%Document Class
\documentclass[a4paper,12pt,reqno]{amsart}


  \headheight0.6in
  \headsep22pt
  \textheight23cm
  \topmargin-1.7cm
  \oddsidemargin 0.5cm
  \evensidemargin0.5cm
  \textwidth15.3cm

%Packages 
\usepackage{amscd,amssymb,amsmath,amsfonts,amsthm, mathrsfs,xspace}
\usepackage{graphicx}
\usepackage{verbatim,enumerate,paralist}
\usepackage{textcomp}
%\usepackage[authoryear]{natbib}
\usepackage[pagebackref,colorlinks]{hyperref}
\hypersetup{%
  urlcolor=blue,linkcolor=blue,citecolor=blue
}
\usepackage{pdfpages}
%\usepackage{pdfsync}%  remove this when finalised

% import these fonts by hand here to avoid clashes with usual blackboard bold usage.
%\pdfmapfile{+bbold.map}
\newcommand{\bbefamily}{\fontencoding{U}\fontfamily{bbold}\selectfont}
\newcommand{\textbbe}[1]{{\bbefamily #1}}
\DeclareMathAlphabet{\mathbbe}{U}{bbold}{m}{n}

% Macro to typeset part of a math formula at a bigger size
\def\makebigger#1#2#3{\hbox{#1$\mathsurround=0pt #2{#3}$}}
\def\bigger#1#2{{\relax\mathpalette{\makebigger#1}{#2}}}

% Macro for the category Delta
\newcommand{\DelCat}{{\bigger\large{\mathbbe{\Delta}}}}
\newcommand{\DelCAT}{{\bigger\Huge{\mathbbe{\Delta}}}}


\usepackage[T1]{fontenc}
\usepackage[all,2cell,poly,knot,tips]{xy}
\UseAllTwocells
% \CompileMatrices

\def\blue{\color{blue}}
\def\red{\color{red}}
\def\black{\color{black}}

%\usepackage{xypdf} % may not be necessary
\setlength{\marginparwidth}{2cm}
\newcommand{\missing}[1]{\marginpar{\color{red}missing:\\#1}}

%Theoretic Environment Setup
\newcounter{Definitioncount}
\setcounter{Definitioncount}{0}

\newtheorem*{Theorem}{Theorem}% no counter
\newtheorem{theorem}{Theorem}[section] % with counter
\newtheorem{lemma}[theorem]{Lemma}
\newtheorem{proposition}[theorem]{Proposition}
% \newtheorem*{Proposition}{Proposition}
\newtheorem{corollary}[theorem]{Corollary}
% \newtheorem*{Corollary}{Corollary}% no counter

\newtheorem*{Conjecture}{Conjecture}% no counter
\newtheorem*{Lemma}{Lemma}% no counter

\theoremstyle{definition}
\newtheorem*{Definition}{Definition}% no counter
\newtheorem*{Observation}{Observation}% no counter
\newtheorem*{Observations}{Observations}% no counter
\newtheorem*{Remark}{Remark}% no counter
\newtheorem{remark}[theorem]{Remark}
\newtheorem*{Example}{Example}% no counter
\newtheorem{example}[theorem]{Example}
\newtheorem{assumption}[theorem]{Assumption}

\renewcommand{\theexample}{\arabic{example}}

\newtheoremstyle{fact}{\bigskipamount}{\medskipamount}{\upshape}{}{\itshape}{. }{ }{Fact}
\theoremstyle{fact}
\newtheorem*{Fact}{Fact}% no counter

\newtheoremstyle{genquest}{\bigskipamount}{\medskipamount}{\upshape}{}{\itshape}{. }{ }{General Question}
\theoremstyle{genquest}
\newtheorem*{GenQuest}{General Question}% no counter

% use \newtheoremstyle to underline the header
\newtheoremstyle{step}{2\bigskipamount}{\medskipamount}{\upshape}{}{\itshape}{. }{ }{\underline{Step~\thestep}}
\theoremstyle{step}
\newtheorem{step}{Step}[subsection]
\renewcommand{\thestep}{\arabic{step}}

\newcommand{\dd}{\colon}
\newcommand{\lra}{\longrightarrow}
\newcommand{\lla}{\longleftarrow}
\newcommand{\ra}{\rightarrow}
\newcommand{\la}{\leftarrow}
\newcommand{\Lra}{\Longrightarrow}
\newcommand{\Lla}{\Longleftarrow}
\newcommand{\ldual}[1]{\mathord{{\let\nolimits\relax\sideset{^\wedge}{}{#1}}}}
\newcommand{\laction}[2]{\mathord{{\let\nolimits\relax\sideset{^{#1}}{}{#2}}}}
\newcommand{\conj}[2]{\mathord{{\let\nolimits\relax\sideset{^{#1}}{}{#2}}}}
\newcommand{\xylongto}[2][2pt]{{\UseTips{}\xymatrix @C+#1{ \ar[r]^{#2}&}}}
\newcommand{\ola}{\overleftarrow}
\newcommand{\ora}{\overrightarrow}

\DeclareMathOperator{\im}{im}

%% Script Shortcuts 
\def\CA{{\mathscr A}}
\def\CB{{\mathscr B}}
\def\CC{{\mathscr C}}
\def\CD{{\mathscr D}}
\def\CE{{\mathscr E}}
\def\CF{{\mathscr F}}
\def\CG{{\mathscr G}}
\def\CH{{\mathscr H}}
\def\CI{{\mathscr I}}
\def\CJ{{\mathscr J}}
\def\CK{{\mathscr K}}
\def\CL{{\mathscr L}}
\def\CM{{\mathscr M}}
\def\CN{{\mathscr N}}
\def\CO{{\mathscr O}}
\def\CP{{\mathscr P}}
\def\CQ{{\mathscr Q}}
\def\CR{{\mathscr R}}
\def\CS{{\mathscr S}}
\def\CT{{\mathscr T}}
\def\CU{{\mathscr U}}
\def\CV{{\mathscr V}}
\def\CW{{\mathscr W}}
\def\CX{{\mathscr X}}
\def\CY{{\mathscr Y}}
\def\CZ{{\mathscr Z}}
\def\ordm{{\mathbf{m}}}
\def\ordn{{\mathbf{n}}}
\def\ordk{{\mathbf{k}}}


\newcommand{\ox}{\otimes}
\newcommand{\x}{\times}

\newcommand{\op}{\ensuremath{^{\textnormal{op}}}}


\def\theequation{{\thesection.\arabic{equation}}}


\begin{document}

\title{Combinatorial categorical equivalences; \\ post arXiv posting}

\author{Stephen Lack}
\address{Department of Mathematics, Macquarie University, NSW 2109, Australia}
\email{steve.lack@mq.edu.au}

\author{Ross Street}
\address{Department of Mathematics, Macquarie University, NSW 2109, Australia}
\email{ross.street@mq.edu.au}

\thanks{Both authors gratefully acknowledge the support of the Australian Research Council Discovery Grant DP130101969; Lack acknowledges with equal gratitude the support of an Australian Research Council Future Fellowship}

\date{\today}
\maketitle

\tableofcontents

\section{Introduction}

Here are some notes on thoughts since submitting \url{arXiv:1402.7151}. 
They concern a general setting with a view to including the example of 
$\Omega_n$ in the sense of \cite{UPMT}.

The intention of this paper is to prove a class of equivalences of categories that 
seem of interest in enumerative combinatorics as per \cite{Species},
representation theory as per \cite{CEF2012}, 
and homotopy theory as per \cite{Berger1996}.
More specifically, for a class of categories $\CP$, we construct a category $\CD$
with zero morphisms (that is, $\CD$ has homs enriched in the category
$1/\mathrm{Set}$ of pointed sets) and an equivalence of categories of the form
\begin{equation}\label{basicequ}
[\CP,\CX] \simeq [\CD,\CX]_{\mathrm{pt}} \ .
\end{equation}
On the left-hand side we have the usual category of functors from
$\CP$ into any additive category $\CX$ which has finite direct sums
and splitting for idempotents. 
On the right-hand side we have the category of functors which preserve
the zero morphisms.

One example has $\CP = \Delta_{\bot,\top}$, the category whose objects are 
finite ordinals with first and last element, and whose morphisms are functions
preserving order and first and last elements. Then $\CD$ is the category with
non-zero and non-identity morphisms
\begin{equation*}
0\stackrel{\partial}\lla 1 \stackrel{\partial}\lla 2\stackrel{\partial}\lla \dots 
\end{equation*}
such that $\partial \circ \partial = 0$.
Since there is an isomorphism of categories 
\begin{equation}\label{interval}
\Delta_{\bot,\top} \cong \Delta_+^{\mathrm{op}} \ ,
\end{equation} 
where the right-hand side is the algebraist's simplicial category (finite ordinals
and all order-preserving functions),
our result \eqref{basicequ} reproduces the Dold-Puppe-Kan Theorem
\cite{Dold, DoldPuppe, Kan}.  

Cubical sets also provide an example of our setting; see Example~\ref{Ex3+}.
We conclude that cubical simplicial abelian groups are equivalent to semi-simplicial
abelian groups. 

For a category $\CA$ equipped with a suitable factorization system 
$(\CE,\CM)$ \cite{FreydKelly},
write $\mathrm{Par}\CA$ (strictly the notation should also show the dependence 
on $\CM$) for the category with the same objects as $\CA$ and
with $\CM$-partial maps as morphisms. We identify $\CE$ with the subcategory of
$\CA$ having the same objects but only the morphisms in $\CE$.
Assume each object of $\CA$ has only finitely many $\CM$-subobjects.
Let $\CX$ be any additive category with finite direct sums and splitting idempotents.
Our main equivalence \eqref{basicequ} includes as a special case an equivalence of categories
\begin{equation}\label{parequiv}
[\mathrm{Par}\CA,\CX] \simeq [\CE,\CX] \ .
\end{equation}
(Here $\CD$ is obtained from $\CE$ by freely adjoining zero morphisms.)
Indeed, we show that this special case implies the general case.

Let $\mathfrak{S}$ be the groupoid of finite sets and bijective functions.
Let $\mathrm{FI}\sharp$ denote the category of finite sets and injective partial functions.
Let $\mathrm{Mod}^R$ denote the category of left modules over the ring $R$.
Our original motivation was to understand and generalize the classification theorem for 
$\mathrm{FI}\sharp$-modules appearing as Theorem 2.24 of \cite{CEF2012},
which provides an equivalence 
\begin{equation*}
[\mathrm{FI}\sharp, \mathrm{Mod}^R] \simeq [\mathfrak{S}, \mathrm{Mod}^R]
\end{equation*}
between the category of functors $\mathrm{FI}\sharp \ra \mathrm{Mod}^R$
and the category of functors $\mathfrak{S} \ra \mathrm{Mod}^R$. 
This is the special case of \eqref{parequiv} above in which $\CA$ is
the category $\mathrm{FI}$ of finite sets and injective functions, 
and $\CM$ consists of all the morphisms. 
This result has provided a new viewpoint on representations of the 
symmetric groups, and a new viewpoint on Joyal species \cite{Species, AnalFunct}.

In order to consider stability properties of representations of the symmetric groups, 
the authors of \cite{CEF2012} also consider $\mathrm{FI}$-modules: 
that is, $R$-module-valued functors from the category $\mathrm{FI}$. 
Each $\mathrm{FI}\sharp$-module
clearly has an underlying $\mathrm{FI}$-module, so their Theorem 2.24
shows how symmetric group representations become $\mathrm{FI}$-modules.
One application they give is a structural version of the Murnaghan Theorem
\cite{Mur38, Mur55}, a problem which has its combinatorial aspects \cite{Sta99}. 
We are reminded of the way in which Mackey functors
\cite{Lindner1976, PanSt} give extra freedom to representation theory.

Another instance of an equivalence of the form \eqref{parequiv} is when
$\CA$ is the category of finite sets with its usual (surjective, injective)-factorization
system. Then $\mathrm{Par}\CA$ is equivalent to the 
category of pointed finite sets,
which is equivalent to Graeme Segal's category $\Gamma$ \cite{Segal}. 
After completing this work, we were alerted to Teimuraz Pirashvili's interesting 
paper \cite{Pirash00} which gives this finite sets example, 
makes the connection with Dold-Puppe-Kan, and discusses stable homotopy of 
$\Gamma$-spaces.  

Consider the basic equivalence \eqref{basicequ}. It is really about Cauchy (or Morita)
equivalence of the free additive category on the ordinary category $\CP$
and the free additive category on the category with zero morphisms $\CD$. 
By the general theory of Cauchy completeness (see \cite{ECC} for example),
to have \eqref{basicequ} for all Cauchy complete additive categories $\CX$, it
suffices to have it when $\CX$ is the category of abelian groups. 
Cauchy completeness amounts to existence of absolute limits (see \cite{ACEC})
and, for additive categories,
amounts to the existence of finite direct sums and splittings for idempotents. 

The strategy of proof of our main equivalence \eqref{basicequ} is as follows.
In Section~\ref{gr} we provide general criteria for deducing a $\CV$-functor
between $\CV$-functor $\CV$-categories is an equivalence given that some other
$\CV$-functor between ``smaller'' $\CV$-functor $\CV$-categories is known to
be an equivalence. 
In Section~\ref{setting} we describe the setting for our main result and then,
in Section~\ref{basex} show that the partial map, Dold-Puppe-Kan and 
cube examples satisfy our assumptions. 
We show in Section~\ref{ker} that the setting of Section~\ref{setting} 
leads to an adjunction between functor categories of a particularly combinatorial kind.
Section~\ref{reducto} reduces the proof that the adjunction is an equivalence to the
particular case of partial maps, then Section~\ref{tpmc} treats the partial map case.

\section{A general result}\label{gr}

Let $\CX$ be a bicategory. 
A morphism $f \dd A\ra B$ is called a {\em map} when it has a right adjoint $f^*\dd B\ra A$.
The map is called {\em Cauchy dense} when $f^*\dd B\ra A$ is monadic. 

Consider the following diagram of two squares, each commuting up to isomorphism.
\begin{eqnarray}\label{twosquares}
\begin{aligned}
\xymatrix{\ar @{} [dr] | {\stackrel{\cong} \Lla}
B \ar[d]_-{m} \ar[r]^-{k^*} & \ar @{} [dr] | {\stackrel{\cong} \Lra} D \ar[d]_-{n} \ar[r]^-{k} & B \ar[d]^-{m}\\
A \ar[r]_-{h^*}           & C  \ar[r]_-{h} & A
}
\end{aligned}
\end{eqnarray}

\begin{proposition}\label{nequivthenmequiv}
In the diagram \eqref{twosquares}, suppose $h$ and $k$ are Cauchy dense maps,
and that the two 2-cells are mates under adjunction in the sense of \cite{KelSt1974}.
If $n$ is an equivalence then so is $m$.  
\end{proposition}
\begin{proof}
By applying the pseudofunctor $\CX(X,-)$, the Proposition reduces to a result in $\mathrm{Cat}$.
Take a monad $S$ on a category $\CC$, a monad $T$ on a category $\CD$,
and a natural isomorphism
$$
\xymatrix{\ar @{} [dr] | {\stackrel{\theta} \Lra}
\CD^T \ar[d]_-{U} \ar[r]^-{M} & \CC^S \ar[d]^-{V} \\
\CD \ar[r]_-{N}           & \CC \ ,
}$$
where $\CC^S$ is the category of Eilenberg-Moore algebras for $S$, similarly $\CD^T$
for $T$,
and $V$ and $U$ are the underlying functors. 
Let $G\dashv V$ and $F\dashv U$ be the usual canonical adjoints.
Let $\phi \dd GN \Rightarrow MF$ be the mate of $\theta$ under adjunction.
We must show that $N$ an equivalence and $\phi$ invertible imply $M$ an equivalence.
It is readily checked that the natural transformation
$$\psi = (SN = VGN \stackrel{V\phi}\lra VMF\stackrel{\theta^{-1}F}\lra NUF = NT)$$
defines a monad morphism 
$$
\xymatrix{\ar @{} [dr] | {\stackrel{\psi} \Lla}
\CD \ar[d]_-{T} \ar[r]^-{N} & \CC \ar[d]^-{S} \\
\CD \ar[r]_-{N}           & \CC \ ,
}$$
inducing a functor $N^{\psi} \dd \CD^T \ra \CC^S$ isomorphic to $M$.
If $\phi$ is invertible, so too is $\psi$. So then, if $N$ is an equivalence, the morphism 
$(N,\psi) \dd (\CD,T) \ra (\CC,S)$ is an equivalence in the 2-category 
$\mathrm{Mnd}(\mathrm{Cat}) = \mathrm{Lax}(1,\mathrm{Cat})$ of monads in $\CX$.
Since the Eilenberg-Moore construction 
$(-)^{(-)}\dd \mathrm{Mnd}(\mathrm{Cat}) \ra \mathrm{Cat}$ is a 2-functor, 
it follows that $N^{\psi}$ is an equivalence and so $M$ is. 
\end{proof}

Let $\CV$ be a symmetric monoidal closed category which is complete and cocomplete.
Suppose $M\dd \CB\ra \CA$ is a $\CV$-module viewed as a
$\CV$-functor $M\dd \CA^{\mathrm{op}} \otimes \CB \ra \CV$.
We call $M$ the {\em kernel} of the adjunction
\begin{equation}\label{keradj} 
\xymatrix @R-3mm {
[\CA,\CV] \ar@<1.5ex>[rr]^{\widehat{M}} \ar@{}[rr]|-{\bot} && [\CB,\CV]
 \ar@<1.5ex>[ll]^{\widetilde{M}} \ ,
}
\end{equation}
where 
\begin{eqnarray}\label{endtilde}
\widetilde{M}(T)A = \int_B{\CV(M(A,B),TB)} = [\CB,\CV](M(A,-),T)
\end{eqnarray}
and
\begin{eqnarray}\label{coendhat}
\widehat{M}(F)B = \int^B{M(A,B)\otimes FA \ .}
\end{eqnarray}

To say $M\dd \CB\ra \CA$ is an equivalence as a $\CV$-module
is the same as saying $\widetilde{M}\dd [\CB,\CV]\ra [\CA,\CV]$
is an equivalence of $\CV$-categories. We also say $\CA$
and $\CB$ are {\em Cauchy} (or {\em Morita}) equivalent when
such an equivalence exists. See \cite{ECC} for more detail.    

A $\CV$-functor $H\dd \CC \ra \CA$ is Cauchy dense 
(as a morphism $H_*=\CA(1_{\CA},H)\dd \CC \ra \CA$
in the bicategory $\CV\text{-}\mathrm{Mod}$) when the canonical
$$\int^C{\CA(HC,A_1)}\otimes \CA(A,HC)\ra \CA(A,A_1)$$ 
is a strong epimorphism in $\CV$ for all $A, A_1\in \CA$. 
When $\CV$ is the category $\mathrm{Agp}$ of abelian groups this means that each
$A\in \CA$ is a retract of a finite direct sum of objects $HC$ in the image of $H$. 

Suppose $M\dd \CB\ra \CA$ and $N\dd \CD\ra \CC$ are $\CV$-modules, and 
$H\dd \CC \ra \CA$ and $K\dd \CD\ra \CB$ are $\CV$-functors.
Each $\CV$-natural family 
$$\theta_{C,D}\dd N(C,D)\ra M(HC,KD)$$ 
induces $\CV$-natural families $\theta^{r}_{A,D}\dd$
$$\int^C{N(C,D)\otimes \CA(A,HC)} \stackrel{\int^C{\theta \otimes 1}}\lra \int^C{M(HC,KD)\otimes \CA(A,HC)} \stackrel{\mathrm{act^r}} \lra M(A,KD)$$ 
and $\theta^{\ell}_{C,B}\dd$
$$\int^D{\CB(KD,B)\otimes N(C,D)} \stackrel{\int^D{1 \otimes \theta}}\lra \int^D{\CB(KD,B)\otimes M(C,D)} \stackrel{\mathrm{act^{\ell}}} \lra M(HC,B) \ .$$

\begin{corollary}\label{gencor}
Suppose $H$ and $K$ are Cauchy dense and both $\theta^{\ell}$ and $\theta^{r}$
are invertible. If $\widetilde{N}\dd [\CD,\CV]\ra [\CC,\CV]$ is an equivalence then $\widetilde{M}\dd [\CB,\CV]\ra [\CA,\CV]$ is too. 
\end{corollary}

\section{The setting}\label{setting}

Let $\CP$ be a category with a subcategory $\CM$ containing all the objects.

Assume we have a functor $(-)^*\dd \CM^{\mathrm{op}} \ra \CP$
which is the identity on objects and is 
such that $m^*\circ m=1$ for all $m\in \CM$.
In particular, the morphisms in $\CM$ are split monomorphisms (coretractions)
in $\CP$. 

We write $\mathrm{Sub}A$ for the (partially) ordered set of isomorphism 
classes of morphisms
$m:U\ra A$ in $\CM$. We will use the term {\em subobject} 
rather than ``$\CM$-subobject'' for these elements. 
The order on $\mathrm{Sub}A$ is the usual one. 
Abusing notation for this order, we simply write $U\preceq V$ when there 
exists an $f : U\ra V$ such that $m=n\circ f$, where $m:U\ra A$ and $n:V\ra A$ in $\CM$. 
There can be confusion when $U=V$ as objects, so we assure the reader that we will take care.
A subobject of $A$ is {\em proper} when it is represented by a non-invertible $m:U\ra A$
in $\CM$; we write $U\prec A$ or $U\prec_m A$. 

Define $\CR$ to be the class of morphisms $r\in \CP$ with the property that, if $r = m\circ x \circ n^*$ with $m, n \in \CM$, then $m, n$ are invertible. 

Define $\CD$ to be the category with zero morphisms (that is, $1/\mathrm{Set}$-enriched
category) obtained from $\CR$ by adjoining zero morphisms.
Composition in $\CD$ of morphisms in $\CR$ is as in $\CP$ if the result is itself
in $\CR$, but zero otherwise. 

Define $\CS$ to be the class of morphisms in $\CP$ of the form $r\circ m^*$ with
$m\in \CM$ and $r\in \CR$. 

Recall that the limit of a diagram consisting of a family of morphisms into
a fixed object $A$ is called a {\em wide pullback}; the morphisms in the
limit cone are called {\em projections}. 
The dual is {\em wide pushout}.  

\begin{assumption}\label{ass1}
Wide pullbacks of families of morphisms in $\CM$ exist, have projections in $\CM$, and become wide pushouts under $m \mapsto m^*$.
\end{assumption}
\begin{assumption}\label{ass2}
If $r, r^{\prime}\in \CR$ are composable then $r^{\prime}\circ r \in \CS$. 
\end{assumption}
\begin{assumption}\label{ass3}
If $r,m^*\circ r \in \CR$ and $m\in \CM$ then $m$ is invertible.
\end{assumption}
\begin{assumption}\label{ass4}
The class $\CM \circ \CM^*$ of morphisms of the form $m\circ n^*$ with $m,n\in \CM$ is closed under composition.
\end{assumption}
\begin{assumption}\label{ass5}
For all objects $A\in \CP$, the ordered set $\mathrm{Sub}A$ is finite. 
\end{assumption}

\begin{proposition}\label{factorize}
Any morphism $f\in \CP$ factors as $f=n\circ r\circ m^*$, uniquely up to isomorphism
for $m,n\in \CM$ and $r\in \CR$. 
\end{proposition}
\begin{proof}
Take $f : A \ra B$ in $\CP$. We use Assumption~\ref{ass1} twice.
Let $n:Y\ra B$ be the wide pullback of all those morphisms 
$V\ra B$ in $\CM$ through which $f$ factors.
Then $n\in \CM$ and there exists a unique $f_1$ with $f = n\circ f_1$.
Let $m:X\ra A$ be the wide pullback of those $U\ra A$ in $\CM$ 
whose left adjoint in $\mathbb{P}$ the morphism $f_1$ factors through.
Then $m\in \CM$ and $f_1 = r \circ m^*$ for a unique $r$.
Clearly $r\in \CR$ and we have uniqueness by a familiar argument. 
\end{proof}

\begin{proposition}\label{oldeff}
If $t\circ s = m\circ r$ with $s,t\in \CS$, $r\in \CR$ and $m\in \CM$ then both $s$ and
$t$ are in $\CR$. 
\end{proposition}
\begin{proof}
First we prove the weaker form: 

{\em if $s\circ r = m\circ r^{\prime}$ with $r,r^{\prime} \in \CR$, $s\in \CS$ and $m\in \CM$ then $s\in \CR$}.

\noindent In obvious notation, put $s=r_1\circ m_1^*$, $m_1^*\circ r = m_2\circ r_2\circ m_3^*$
and $r_1\circ m_2=m_4\circ r_3\circ m_5^*$.
Then $m\circ r^{\prime} = s\circ r = r_1\circ m_1^*\circ r = r_1\circ m_2\circ r_2 \circ m_3^* = m_4\circ r_3\circ m_5^*\circ r_2 \circ m_3^*$. 
Let $p,q\in \CM$ be the projections in the pullback of $m,m_4$. 
Then there exists a unique $u$ into the pullback with $r^{\prime} = p\circ u$
and $r_3\circ m_5^*\circ r_2 \circ m_3^* = q\circ u$. 
Since $r^{\prime}$ can factor through no proper $\CM$, we deduce that $p$ is invertible.
Put $n = q\circ p^{-1}$.
Then $m=m_4\circ n$ and $n\circ r^{\prime} = q\circ u = r_3\circ m_5^*\circ r_2 \circ m_3^*$. Therefore $r^{\prime} = n^* \circ r_3\circ m_5^*\circ r_2 \circ m_3^*$;
by definition of $\CR$, we obtain that $m_3$ is invertible. 
Then $m_1^*\circ (r\circ m_3) = m_2 \circ r_2$ implies 
$(m_1\circ m_2)^*\circ (r\circ m_3) = r_2$. Assumption~\ref{ass3} applies to yield that
$m_1\circ m_2$ is invertible. So $m_1$ has a right inverse, as well as its left inverse
$m_1^*$, and so is invertible. So $s=r_1\circ m_1^*$ is in $\CR$ as asserted.  

Now we come to the proof of the Proposition. Since $s\in \CS$, $s= r_1\circ m_1^*$. 
Then, by definition of $\CR$, $m^*\circ t \circ r_1\circ m_1^* 
= m^*\circ t \circ s= m^*\circ m\circ r = r$ implies $m_1$
invertible. So $s\in \CR$. Now the weaker form above applies to yield $t\in \CR$.  
\end{proof}

For each $u\dd A\ra B$ in $\CP$, put $u=m_u\circ s_u$ where $m_u\in \CM$
and $s_u\in \CS$; we are using Proposition~\ref{factorize}.

\begin{proposition}\label{esssub}
Suppose $u\dd A\ra B$ and $v\dd B\ra C$ in $\CP$. 
\begin{enumerate}
\item[(a)] If $s_{vu}\in \CR$ then $s_u\in \CR$. 
\item[(b)] If $s_{vu}\in \CR$ and $u\in \CS$ then $u, s_v\in \CR$.
\end{enumerate}
\end{proposition}
\begin{proof}
We prove (b) first. We have $v\circ u = m_v\circ s_v \circ u = m_v \circ m_{s_vu} \circ s_{s_vu}$. By uniqueness of factorization, we may take $s_{vu}= s_{s_vu}$.
So $s_v\circ u = m_{s_vu} \circ s_{vu}$. By Proposition~\ref{oldeff}, if $s_{vu}\in \CR$ and $u\in \CS$ then $s_v, u\in \CR$. 
Now we prove (a). We have $v\circ u = v \circ m_u\circ s_u$.
By (b), $s_{vm_us_u} = s_{vu} \in \CR$ implies $s_u\in \CR$.  \end{proof}

\section{Basic examples}\label{basex}

\begin{example}\label{Ex1+}
Take a category $\CA$ with a factorization system $(\CE,\CM)$ in the sense of \cite{FreydKelly}. 
Assume that the pullback of any morphism with one in $\CM$ exists. 
Assume every morphism in $\CM$ is a monomorphism and every object
of $\CA$ has only finitely many $\CM$-subobjects. 
A span $f = (X\stackrel{f_0}\lla U \stackrel{f_1}\lra Y)$ is called a {\em partial map}
$f : X\lra Y$ when $f_0$ is in $\CM$. 
If $g = (X\stackrel{g_0}\lla V \stackrel{g_1}\lra Y)$ is another partial map,
we write $f\le g$ when there exists $h:U\ra V$ with $g_0\circ h=f_0$ and $g_1\circ h=f_1$.  
Let $\CP= \mathrm{Par}\CA$ 
denote the category
whose objects are all those of $\CA$ and whose morphisms are isomorphism
classes $[f]$ of partial maps. Composition is that of spans: that is, by pullback.
Taking the usual order $f\le g$ on partial maps, we can regard $\mathrm{Par}\CA$
as a locally partially ordered 2-category.
We identify $f\dd X \ra Y$ in $\CA$ with the morphism $[1_X,X,f] \dd X \lra Y$ in $\CP$.
In this way, we have the $\CM$ we require for $\CP$ as the one in $\CA$.
For $m\dd U\ra X$ in $\CM$, a right adjoint in $\mathrm{Par}\CA$ is defined by 
$m^*= [m,U,1_U] : X\ra U$, and clearly $m^*\circ m = 1$.
This gives our functor $(-)^*\dd \CM^{\mathrm{op}} \ra \CP$. 
Every partial map $f = (X\stackrel{f_0}\lla U \stackrel{f_1}\lra Y)$ has 
$$[f] = f_{1}\circ f_0^* $$
where $f_0\in \CM$. Furthermore, $f_1 = m\circ e$ uniquely up to isomorphism
for $m\in \CM$ and $e\in \CE$. 
It follows therefore that $\CR=\CE$ and that $[f]\in \CS$ if and only if $f_1\in \CE$.

Now we look at our Assumptions. To see that Assumption~\ref{ass1} holds, first note
that since we are assuming the $\CM$-subobjects form a finite set, 
finite wide pullbacks can be obtained from pullbacks.
By assumption, pullbacks of $\CM$s exist in $\CA$ and a pullback of an 
$\CM$ is an $\CM$ in a factorization system.
It is a pleasant exercise to see that these pullbacks remain pullbacks in $\CP$
and become pushouts in $\CP$ on taking left adjoints. 

We certainly have Assumption~\ref{ass2}; 
indeed $\CR$ is closed under composition because $\CE$ is.

For Assumption~\ref{ass3}, take $r=[1_A,A,e]$ with $e\in \CE$ and 
$m=[1_V,V,m]$ in $\CM$. To have $m^*\circ r = [n,P,u] \in \CR$, we must have
$n$ invertible and $u\in \CE$. Then $e=m\circ u \circ n^{-1}$ implies $m\in \CE$;
so $m$ is invertible. 

The class of partial maps in Assumption~\ref{ass4} are those
of the form $[m,U,n]$ with $m,n\in \CM$; these are closed under composition since
pullbacks of $\CM$s exist and are in $\CM$. 

Assumption~\ref{ass5} was one of our assumptions on the factorization system on $\CA$. 

Notice that $\CS$ is not closed under composition unless each pullback of an $\CE$
along an $\CM$ is an $\CE$. This is true in many examples. 
\end{example}

\begin{example}\label{Ex2+}
Take $\CP$ to be the category $\Delta_{\bot,\top}$ of finite
non-empty ordinals $n = \{0,1,\dots,n-1\}$ with morphisms those functions which preserve
first element, last element and order.
Functors out of this category are augmented simplicial objects because of the isomorphism \eqref{interval}. 
Morphisms $\xi, \zeta : m\ra n$ are ordered by taking $\xi \le \zeta$ iff and only if 
$\xi(i) \le \zeta(i)$ for all $i\in m$. This makes $\Delta_{\bot,\top}$ into a locally
partially ordered 2-category;
it is a locally full sub-2-category of $\mathrm{Cat}$. 
Take $\CM$ to consist of all the injective functions in $\Delta_{\bot,\top}$;
each such injection $\partial$ has a left adjoint $\partial^*$ 
(and also a right adjoint for that matter) in $\Delta_{\bot,\top}$; clearly $\partial^*$ is
surjective with $\partial^* \circ \partial = 1$ since $\partial$ is a fully faithful functor.
A surjection $\sigma$ in $\CP$ is of the form $\partial^*$ if and only if $\sigma(i)=0$
implies $i=0$.    
We write $\sigma_k : m+1 \ra m$ for the order-preserving surjection which takes the value $k$ twice.
We write $\partial_k : m \ra m+1$ for the order-preserving injection which does not have $k$ in its image. 
Note that $\partial_0\notin \CP$, while $\sigma_k \dashv \partial_k  \dashv \sigma_{k-1}$ 
in $\mathbb{P}$ and $\sigma_k$ is a $\partial^*$ if and only if $k>0$.
Every $\xi \in \CP$ factors uniquely as 
$$\xi = \partial_{i_s}\circ \dots \circ \partial_{i_1}\circ \sigma_{j_1}\circ \dots \circ \sigma_{j_t}$$ 
for $0< i_1<\dots < i_s < n-1$ and $0\le j_1<\dots < j_t < m-1$. 

We claim $\CR$ consists of the identities and the surjections $\sigma_0 : m+1 \ra m$.
The only invertible morphisms in $\CP$ are identities.
Since members of $\CR$ factor through no proper injection, they must be surjective.
Every surjection $\tau$ is either of the form $\partial^*$ or uniquely of the form 
$\sigma_0\circ \partial^*$. 
Neither of these forms is permissible for $\tau \in \CR$ unless the injection $\partial$
is an identity. This proves our claim. 
It is also clear then that $\CS$ consists of all the surjections in $\CP$. 

To prove the Assumptions, we make use of the simplicial identities 
(see page 24 of \cite{GabZis} for example) which, 
apart from $\sigma_i\circ \partial_i = 1 = \sigma_{i-1}\circ \partial_i$,
say that, for all $i<j$, 
$$\partial_j\circ \partial_i = \partial_i \circ \partial_{j-1} \ , \ \sigma_{j-1}\circ \sigma_i = \sigma_i \circ \sigma_j \ , \ \sigma_{j}\circ \partial_i = \partial_i \circ \sigma_{j-1} \ , \ \sigma_{i}\circ \partial_{j+1} = \partial_j \circ \sigma_i \ .$$

Assumption~\ref{ass1} follows from the fact that the pullback of any two 
monomorphisms in $\CP$ exists and is absolute (that is, preserved by all functors); 
see page 27 of \cite{GabZis} on the Eilenberg-Zilber Theorem. 
Alternatively, notice that the squares
$$
\xymatrix{
n-1 \ar[rr]^-{\partial_{j-1}} \ar[d]_-{\partial_{i}} && n \ar[d]^-{\partial_{i}} \\
n \ar[rr]_-{\partial_{j}}  && n+1 }
\qquad
\xymatrix{
n-1  && n \ar[ll]_-{\sigma_{j-1}} \\
n \ar[u]^-{\sigma_i}  && n+1 \ar[u]_-{\sigma_i} \ar[ll]^-{\sigma_j}}
$$
are respectively a pullback and pushout in $\CP$ for $0< i < j < n$.
Then the general result follows by stacking these squares vertically and horizontally.

In the present example, $\CS$ is closed under composition, making Assumption~\ref{ass2} clear.

For Assumption~\ref{ass3}, suppose we have $\partial \circ \rho \in \CR$ and $\rho \in \CR$. Since $\rho$ is surjective, $\partial^* \circ \rho = 1$ implies $\rho = 1$ 
and hence $\partial = 1$, as required. Otherwise $\partial^* \circ \rho = \sigma_0$.
Yet, if $\rho = 1$ this contradicts the lack of right adjoint for $\sigma_0$.
So $\rho = \sigma_0$. Thus $\partial^* \circ \sigma_0 = \sigma_0$, and we can
cancel $\sigma_0$ to obtain again $\partial = 1$. 

For Assumption~\ref{ass4}, the class $\CM\circ \CM^*$ of morphisms consists of
those which reflect $0$. That is, $\xi = \mu \circ \partial^*$ with
$\mu\in \CM$ if and only if $\xi(i)=0$ implies $i=0$. This class is
clearly closed under composition. 

Finally, Assumption~\ref{ass5} is clear. 

Notice that the arguments above equally apply to the full subcategory  
$\Delta_{\bot\neq\top}$ of $\Delta_{\bot,\top}$ obtained by removing the object $1$.
Functors $\Delta_{\bot\neq \top}\ra \CX$ are the traditional simplicial objects in $\CX$. 
\end{example}

\begin{example}\label{Ex3+} This example is about the cubical category $\mathbb{I}$ as used by Sjoerd Crans \cite{CransThesis} and Dominic Verity \cite{LDTHOC, Verity07}. 
Functors with domain $\mathbb{I}$ are cubical objects in the codomain category. 
Verity constructed $\mathbb{I}$ as the free monoidal category containing a cointerval.

For each natural number $k$, define a poset $\langle k\rangle = \{-,1,2,\dots ,k,+\}$ 
by adjoining a bottom element $-$ and a top element $+$ to the discrete poset $\{1,2,\dots ,k\}$.
Any function $f : \langle k\rangle \ra \langle h\rangle$ which preserves top and bottom is order-preserving.
Thus we get a locally partially ordered 2-category with objects the $\langle k\rangle$, 
with morphisms the top-and-bottom-preserving functions, and with the pointwise order. 
Let $\mathbb{I}$ be the locally full sub-2-category
consisting of those $f:\langle k\rangle \ra \langle h\rangle$ for which,
if $f(i),f(j)\notin \{-,+\}$ then $i<j$ if and only if $f(i)<f(j)$. 

Let $\CP$ be the underlying category of this $\mathbb{I}$.
Let $\CM$ consist of the morphisms in $\CP$ which are injective as functions. 
Given such an $m : \langle k\rangle \ra \langle h\rangle$ in $\CM$,
define $m^* : \langle h\rangle \ra \langle k\rangle$ to send each $m(i)$
in the image of $m$ to $i$ and everything else to $+$. 
Clearly $m^*\in \CP$ and $m^*\circ m = 1$.
Furthermore, $mm^*(j)$ is equal to $j$ if $j=m(i)$ for some $i$, and $+$ otherwise.
Therefore $1\le m\circ m^*$ showing $m^*$ to be left adjoint to $m$
with identity counit.

We can characterize morphisms of the form $m^*$ as those which are surjective as functions and reflect the bottom element $-$.
Consequently $\CR$ consists of the morphisms which are surjective as
functions and reflect the top element $+$.

Assumption~\ref{ass5}~and~\ref{ass2} are clear. 
\begin{equation}\label{cubpbpo}
\begin{aligned}
\xymatrix{
\langle \ell \rangle \ar[rr]^-{q} \ar[d]_-{p} &&  \langle v \rangle \ar[d]^-{n} \\
\langle u \rangle \ar[rr]_-{m}  && \langle k \rangle}
\qquad
\xymatrix{
\langle k \rangle \ar[rr]^-{n^*} \ar[d]_-{m^*} && \langle v \rangle \ar[d]_-{q^*} \ar@/^/[ddr]^g \\
\langle u \rangle \ar[rr]^-{p^*} \ar@/_/[drrr]_{f} && \langle \ell \rangle \\
&& & \langle h \rangle }
\end{aligned}
\end{equation}

For Assumption~\ref{ass1}, the existence of intersections is obvious.
However, we must show that taking left adjoints gives cointersections.
Take a pullback as in the left-hand diagram of \eqref{cubpbpo} 
and consider the right-hand diagram.
Assume $f\circ m^*=g\circ n^*$.
It suffices to show $f\circ p\circ p^*=f$.
Now $fpp^*(i)$ is equal to $fp(j)$ if $i=p(j)$ for some $j$, and equal to $+$ otherwise.
In the first case, we have $fpp^*(i)=fp(j)=f(i)$, as required.
In the second case, if $i$ does not have the form $p(j)$ then
$m(i)$ is not in the image of $n$, and so $gn^* m(i)=+$;
thus $f(i)=fm^*m(i)=+$ and we again have $fpp^*(i)=f(i)$.

For Assumption~\ref{ass3}, suppose $m\in \CM$ and $r, m^*\circ r\in \CR$. 
Suppose $m^*(i)=+$. 
Since $r$ is surjective, we have $i=r(j)$ so that $m^*r(j)=+$.
However $m^*\circ r$ reflects $+$. So $j=+$, yielding $i=r(+)=+$.
This proves $m^*$ reflects $+$ and therefore must be invertible.

For Assumption~\ref{ass4}, the composites of the form $m\circ n^*$
are clearly the (not necessarily surjective) morphisms which reflect $-$.
Clearly these are closed under composition.

There is another possible characterization of $\CR$, namely as the category of right adjoints to the morphisms in $\CM$. 
Thus in fact $\CR$ is dual to $\CM$. 
Now $\CM$ is really just the category $\Delta_{\mathrm{inj}}$ of finite ordinals and injective order-preserving maps, 
and $\CR \cong \CM^{\mathrm{op}}$.
Also the category $\CM^*$, with morphisms the $m^*\in \CM$, is dual to $\CM$. 
The factorization of Proposition~\ref{factorize} in this case shows the category $\CP$ is a composite 
$$\mathbb{I} = \Delta_{\mathrm{inj}}\circ \Delta_{\mathrm{inj}}^{\mathrm{op}}\circ \Delta_{\mathrm{inj}}^{\mathrm{op}} \ ,$$
relative to suitably defined distributive laws.

An alternative viewpoint is that $\mathbb{I}$ is 
$\mathrm{Par}\mathrm{Par}\Delta_{\mathrm{inj}}$,
where in each case partial maps are defined relative to the morphisms in 
$\Delta_{\mathrm{inj}}$. 
\end{example}

\section{The kernel module}\label{ker}

In the setting of Section~\ref{setting}, we shall define a functor
\begin{eqnarray}\label{moduleM}
M\dd \CD^{\mathrm{op}}\times \CP \lra \mathrm{Agp}
\end{eqnarray}
which preserves zeros in the first variable.

Let $\mathbb{Z}\Lambda$ denote the free abelian group on the set $\Lambda$.
This defines the value on objects of a strong monoidal functor 
$\mathbb{Z} \dd \mathrm{Set} \ra \mathrm{Agp}$
whose right adjoint is the forgetful functor.

For each $u\dd A\ra B$ in $\CP$, we have $m_u\in \CM$ and $s_u\in \CS$ as in the triangle
\begin{eqnarray}\label{fact_u}
\begin{aligned}
\xymatrix{
A \ar[rd]_{s_u}\ar[rr]^{u}   && B  \\
& S_u \ar[ru]_{m_u}  &
} 
\end{aligned}
\end{eqnarray}  
by using Proposition~\ref{factorize}.

Define $M$ on objects by
\begin{eqnarray}\label{Mobj}
M(A,B) = \mathbb{Z}\{u\in \CP(A,B) \dd s_u\in \CR \} \ .
\end{eqnarray}
Suppose $r\dd A_1 \ra A$ in $\CR$ and $f\dd B\ra B_1$ in $\CP$.
Define $M$ on morphisms by
\begin{eqnarray}\label{Mmor}
M(r,f)u = 
\begin{cases}
f\circ u \circ r & \text{for }  s_{fur}\in \CR \\
0 & \text{otherwise} \ .
\end{cases}
\end{eqnarray}
To verify that this is a functor we need to see that 
$$M(r_1,f_1)M(r,f)u = M(r\circ_{\CD} r_1,f_1\circ f)u$$ holds. 
This amounts to showing
that the left-hand side is non-zero if and only if the right-hand side is, 
because then both sides
equal $f_1\circ f\circ u \circ r\circ r_1$. 
In other words, we need to see that 
$s_{fur}, s_{f_1furr_1}\in \CR$ if and only if $r\circ r_1, s_{f_1furr_1}\in \CR$.
In fact, $s_{f_1furr_1}\in \CR$ implies both $s_{fur}\in \CR$ and $r\circ r_1\in \CR$.   
For we know by Assumption~\ref{ass2} that $r\circ r_1 \in \CS$; 
so, by Proposition~\ref{esssub}(b), $s_{f_1furr_1}\in \CR$ implies $r\circ r_1 \in \CR$.
By Proposition~\ref{esssub}(a), we also conclude that $s_{furr_1}\in \CR$;
then Proposition~\ref{esssub}(b) gives $s_{fur}\in \CR$.    

Let $\mathbb{Z}_{\dagger}\CD$ denote the free $\mathrm{Agp}$-enriched category on
the pointed-set-enriched category $\CD$. Let $\mathbb{Z}_{*}\CP$ denote the free 
$\mathrm{Agp}$-enriched category on the category $\CP$. For any additive (that is,
$\mathrm{Agp}$-enriched) category $\CX$, we have isomorphisms of categories
$$[\mathbb{Z}_{\dagger}\CD,\CX]_{\mathrm{add}} \cong [\CD,\CX]_{\mathrm{pt}}
\ \text{ and } \ [\mathbb{Z}_{*}\CP,\CX]_{\mathrm{add}} \cong [\CD,\CX] \ , $$ 
where, for emphasis, we write the subscript ``$\mathrm{pt}$'' for the pointed-set-enriched
functor category and ``$\mathrm{add}$'' for the $\mathrm{Agp}$-enriched functor category.    

We identify \eqref{moduleM}
with its obvious extension to an additive functor
\begin{eqnarray}\label{additiveM}
M : \mathbb{Z}_{\dagger}\CD^{\mathrm{op}} \otimes \mathbb{Z}_{*}\CP \lra \mathrm{Agp} \ .
\end{eqnarray}
Now we are in the situation for a kernel adjunction of the form \eqref{keradj} with 
$\CV = \mathrm{Agp}$. We have an adjunction.
\begin{equation}\label{settingkeradj} 
\xymatrix @R-3mm {
[\CD,\mathrm{Agp}]_{\mathrm{pt}} \ar@<1.5ex>[rr]^{\widehat{M}} \ar@{}[rr]|-{\bot} && [\CP,\mathrm{Agp}]
 \ar@<1.5ex>[ll]^{\widetilde{M}} \ .
}
\end{equation}
More combinatorial expressions than \eqref{coendhat} and \eqref{endtilde} exist for $\widehat{M}$ and $\widetilde{M}$ in our present situation.
\begin{theorem}\label{combadj}
Suppose $M$ is the kernel as in \eqref{additiveM}. 
For any zero preserving functor $F\dd \CD\ra \mathrm{Agp}$ and any functor 
$T\dd \CP \ra\mathrm{Agp}$, there are bijections
\begin{eqnarray*}
\widehat{M}(F)B \cong \sum_{S\preceq_m B}{FS} \ ,
\qquad
\widetilde{M}(T)A \cong \bigcap_{U\prec_m A} {\mathrm{ker} \ {T(m^*:A\ra U)}} \ .
\end{eqnarray*}
\end{theorem}
\begin{proof}
We shall prove $$\int^{A\in \mathbb{Z}_{\dagger}\CD}{M(A,B)\otimes FA} \cong \sum_{S\preceq_m B}{FS}$$
by showing that the family of morphisms
$$ \chi_A \dd M(A,B)\otimes FA \lra \sum_{S\preceq_m B}{FS} $$
defined by $\chi_A(u\otimes x) = (Fs_u)x\in FS_u$ (using the notation of \eqref{fact_u}) 
is a universal extraordinary natural transformation (that is, a universal wedge). 
In order for a family $\theta_A \dd M(A,B)\otimes FA \ra X$ in $\mathrm{Agp}$ to be
extraordinary natural, we require 
\begin{eqnarray}\label{wedge}
\theta_A(u\otimes (Fr)y) = 
\begin{cases}
\theta_{A_1}((u\circ r)\otimes y) & \text{for }  s_{ur}\in \CR \\
0 & \text{otherwise} \ .
\end{cases}
\end{eqnarray}
for all $r\dd A_1\ra A$ in $\CR$, $u\dd A\ra B$ in $\CP$ with $s_u \in \CR$, and $y\in FA_1$.
Notice that in this situation $s_{ur} = s_u\circ r$ by uniqueness of factorization and Assumption~\ref{ass2}. 
Now for the family $\chi$ we have $\chi_A(u\otimes (Fr)y) = F(s_u\circ_{\CD}r)y$. 
If $s_{ur}\in \CR$ then $s_u\circ_{\CD}r = s_u\circ r = s_{ur}$, so
$F(s_u\circ_{\CD}r)y = F(s_{ur})y = \chi_{A_1}((u\circ r)\otimes y)$.
Otherwise, $F(s_u\circ_{\CD}r) = F(0) = 0$, as required.
For universality, we must see that a general wedge $\theta$ factors uniquely through $\chi$.
Define $\phi \dd \sum_{S\preceq_m B}{FS} \ra X$ by $\phi(x\in FS) = \theta_S(m\otimes x)$.
With $s_u\in \CR$, we have, using \eqref{wedge}, 
$$\theta_A(u\otimes x) = \theta_A((m_u\circ s_u)\otimes x) = \theta_S(m_u\otimes (Fs_u)x) = \phi((Fs_u)x\in FS_u) = \phi \chi_A(u\otimes x) \ .$$
Thus $\theta =  \phi \chi$. By taking $u=m\in \CM$, we also see that the definition of $\phi$ is forced. 

Define $\widetilde{T}A = \bigcap_{U\prec_m A} {\mathrm{ker} \ {T(m^*:A\ra U)}}$,
and, for $r\dd A\ra A_1$, define $(\widetilde{T}r)a = (Tr)a$.
For this we need to see that $a\in \widetilde{T}A$ implies $(Tr)a\in \widetilde{T}A_1$.
Take a non-invertible $n\dd V\ra A_1$. Using Proposition~\ref{factorize}, we have $n^*\circ r = \ell \circ r_1 \circ m^*$. So $(n\circ \ell)^*\circ r = r_1\circ m^*$. If $m$ is invertible then $(n\circ \ell)^*\circ r \in \CR$. By Assumption~\ref{ass3}, $n\circ \ell$ is invertible. So $n$ is invertible, a contradiction. So $m$ is not invertible and we have 
$(Tn^*)(Tr)a = T(\ell)T(r_1)(Tm^*)a = T(\ell)T(r_1)0 = 0$. So $(Tr)a \in \widetilde{T}A_1$.         

Put $\widehat{F}B = \sum_{S\preceq_m B}{FS}$ and transport the functoriality in $B$
across the isomorphism with $\widehat{M}(F)B$.
We will prove that there is a natural isomorphism
$$[\CP,\CX](\widehat{F}, T) \cong [\CD,\CX]_{\mathrm{pt}}(F,\widetilde{T})$$
and hence conclude that $\widetilde{M}(T) \cong \widetilde{T}$.
Take a natural transformation $\theta \dd \widehat{F}\Rightarrow T \dd \CP \ra \CX$. 
The definition of $\widehat{F}f$ involves writing $f\circ m = m_{fm} \circ s_{fm}$ then, 
with $s_{fm}\in \CR$, we have commutativity in the following diagram. 
\begin{equation}\label{thetanat}
\begin{aligned}
\xymatrix{
FU \ar[r]^-{\mathrm{in}_U} \ar[d]_-{Fs_{fm}} & \widehat{F}A \ar[d]^-{\widehat{F}f} \ar[r]^-{\theta_A} & TA \ar[d]^-{Tf}\\
FV \ar[r]_-{\mathrm{in}_V} & \widehat{F}B \ar[r]_-{\theta_B} & TB}
\end{aligned}
\end{equation}
Consider a noninvertible $\ell : W \ra A$ in $\CM$. 
Then, by the properties of $\CR$, $\ell^* \in \CS$ and $\ell^* \notin \CR$.
So $T\ell^*\circ \theta_A \circ \mathrm{in}_{A} = \theta_W \circ \widehat{F}\ell^*\circ \mathrm{in}_{A} = \theta_W \circ 0 = 0$. This implies there exists a unique morphism $\phi_A : FA \lra \widetilde{T}A$ 
such that the following square commutes.
\begin{equation}\label{defphi}
\begin{aligned}
\xymatrix{
FA \ar[rr]^-{\phi_A} \ar[d]_-{\mathrm{in}_A} && \widetilde{T}A \ar[d]^-{i_A} \\
\widehat{F}A \ar[rr]_-{\theta_A} && TA}
\end{aligned}
\end{equation}
Naturality of $\phi$ is proved as follows using \eqref{thetanat} with $f=r \in \CR$:
$i_B\circ \phi_B\circ Fr = \theta_B \circ \mathrm{in}_{B} \circ Fr = \theta_B \circ \widehat{F}r \circ \mathrm{in}_{A} = Tr \circ \theta_A \circ \mathrm{in}_{A} = Tr \circ i_A \circ \phi_A = i_B\circ \widetilde{T}r\circ \phi_A $. 

For the inverse direction, take any natural transformation $\phi \dd F\Rightarrow \widetilde{T} \dd \CD\ra \CX$.
Define $\theta$ by commutativity of the following diagram.
\begin{equation}\label{thetadef}
\begin{aligned}
\xymatrix{
FU \ar[r]^-{\mathrm{in}_U} \ar[d]_-{\phi_U} & \widehat{F}A  \ar[r]^-{\theta_A} & TA \\
\widetilde{T}U \ar[rr]_-{i_U} &  & TU \ar[u]_-{Tm}}
\end{aligned}
\end{equation}       
We need to prove the right-hand square of \eqref{thetanat} commutes when precomposed
with any $\mathrm{in}_U$ for any $m:U\ra A$ in $\CM$.  
In the case where $s_{fm}\in \CR$, the desired commutativity is a consequence of the commutativity of the following three squares.
\begin{equation}\label{pfthetanat}
\begin{aligned}
\xymatrix{
FU \ar[r]^-{\phi_U} \ar[d]_-{Fs_{fm}} & \widetilde{T}U \ar[d]_-{\widetilde{T}s_{fm}} \ar[r]^-{i_A} & TU \ar[r]^-{Tm} \ar[d]^-{Ts_{fm}} & TA \ar[d]^-{Tf}\\
FV \ar[r]_-{\phi_V} & \widetilde{T}V  \ar[r]_-{i_B} & TV \ar[r]_-{Tn_{fm}} & TB}
\end{aligned}
\end{equation}
In the case where $s_{fm}\notin \CR$, we can write $s_{fm} = r\circ \ell^*$ for some noninvertible $\ell \in \CM$.
Then (using both $U\preceq A$ and $U\preceq U$) we have $Tf \circ \theta_A \circ \mathrm{in}_U = T(f\circ m)\circ i_U \circ \phi_U =T(n_{fm}\circ r) \circ T\ell^*\circ i_U \circ \phi_U = T(n_{fm}\circ r) \circ 0 \circ \phi_U = 0 = \theta_B \circ \widehat{F}f \circ \mathrm{in}_U$, as required.

To show that the assignments are mutually inverse, take $\theta$ and define $\phi$ by \eqref{defphi}. 
Let $\bar{\theta}$ be as $\theta$ is in \eqref{thetadef}. Then $\bar{\theta}_A\circ \mathrm{in}_U = Tm \circ i_U\circ \phi_U = Tm \circ \theta_U \circ \mathrm{in}_{U} = \theta_A \circ \widehat{F}m\circ \mathrm{in}_{U} = \theta_A \circ \mathrm{in}_U$. So $\bar{\theta} = \theta$. 

On the other hand, take $\phi$ and define $\theta$ by \eqref{thetadef}. 
Let $\phi^{\prime}$ be as $\phi$ is in \eqref{defphi}. 
Then $i_A \circ \phi^{\prime}_A = \theta_A \circ \mathrm{in}_{A} = i_A\circ \phi_A$. 
So $\phi^{\prime} = \phi$.    
\end{proof}

\section{Reduction of the problem}\label{reducto}

The problem in question is to show the adjunction \eqref{settingkeradj} is an equivalence. 
We wish to use the theory in Section~\ref{gr} to reduce the problem to a special 
case of Example~\ref{Ex1+} involving partial maps. In Section~\ref{tpmc}, we
will treat the general case of Example~\ref{Ex1+}. 

Consider $\CP$, $\CM$, $\CS$, $\CR$ and $\CD$ as in the setting of Section~\ref{setting}.
Let $\CI$ be the class of invertible morphisms in $\CP$.
Let $\CK = \CM\circ \CM^*$ be the subcategory of $\CP$ as assured by 
Assumption~\ref{ass4}.
Then $(\CI,\CM)$ is a factorization system on $\CM$ becoming an instance of 
Example~\ref{Ex1+} with $\mathrm{Par}\CM = \CK$.
Denote the kernel module for this situation by 
$N \dd \CI^{\mathrm{op}} \times \CK \ra \mathrm{Agp}$;
on objects it is given by $N(C,D) = \mathbb{Z}\CM(C,D)$.  

Let $H\dd \CI \ra \CD$ and $K\dd \CK \ra \CP$ be the inclusions;
they are the identity on objects and so Cauchy dense. 

We have the obvious inclusion  
\begin{eqnarray}\label{specialtheta}
\theta_{C,D} \dd N(C,D) \ra M(HC,KD)
\end{eqnarray}
taking $m\in \CM(C,D)$ to $m \in \CP(C,D)$ (using $s_m = 1$).

\begin{proposition}\label{twocoends}
For $\theta \dd N \Rightarrow M(H,K)$ defined at \eqref{specialtheta}, 
the corresponding natural families
$$\theta^r_{A,D} \dd \int^{C\in \CI}{N(C,D)\otimes \CD(A,HC)} \lra M(A,KD) \ ,$$
$$\theta^{\ell}_{C,B} \dd \int^{D\in \CK}{\CP(KD,B)\centerdot N(C,D)} \lra M(HC,B) \ ,$$
as in Section~\ref{gr}, are both invertible. 
\end{proposition}
\begin{proof}
For $r\dd A\ra C$ in $\CR$ and $m\dd C\ra D$ in $\CM$, 
we have $\theta^r_{A,D}[m\otimes r]= m\circ r$. 
Yet every $u\dd A\ra D$ with $s_u\in \CR$ factors uniquely up
to a morphism in $\CI$ as $u=m_u\circ s_u = \theta^r_{A,D}[m_u\otimes s_u]$.
So we have our inverse bijection.

By the definition of the second coend, for all $k\dd D\ra D_1$ in $\CK$ and 
$g\dd D_1\ra B$, we have $[(g\circ k)\otimes m] = [(g\otimes k)\circ m]$
when $k\circ m \in \CM$, zero otherwise.
For $m\dd C\ra D$ in $\CM$ and $f\dd D\ra B$ in $\CP$, we have
$[f\otimes m] = [(f\circ m)\otimes 1] = [u\otimes 1]$ for some $u \in \CP$.
But $u = \ell \circ r \circ n^*$ for some $n,\ell \in \CM$ and $r\in \CR$.
So $[u\otimes 1] = [(\ell \circ r)\otimes n^*]$ provided $n^*\in \CM$, zero
otherwise. However, $n^*\in \CM$ implies $n$ invertible. 
So $s_u\in \CR$. In other words, $[f\otimes m]$ is either zero or of the
form $[u\otimes 1]$ for a unique $u$ with $s_u\in \CR$.  
By definition, $\theta^{\ell}_{C,B}[f\otimes m]= f\circ m$ when $s_{fm}\in \CR$,
zero otherwise. 
We have shown that an inverse to $\theta^{\ell}_{C,B}$ is provided by 
$u \mapsto [u\otimes 1_C]$.
\end{proof}

Combining Proposition~\ref{twocoends} and Corollary~\ref{gencor}, we have:
 
\begin{corollary}
If $\widetilde{N}$ is an equivalence then so is $\widetilde{M}$.
\end{corollary}

\section{A treatment of the partial map case}\label{tpmc}

We begin with any category $\CA$ equipped with a 
factorization system $(\CE,\CM)$ as in Example~\ref{Ex1+}.
We have $\CP = \mathrm{Par}\CA$ and $\CR = \CE$.
Also $\CD$ is $\CE$ with adjoined zero morphisms, so 
$[\CD,\CX]_{\mathrm{pt}} \cong [\CE,\CX]$ for any category
$\CX$ with zero morphisms.  

For every pullback
\begin{equation}\label{pbm}
\xymatrix{
P \ar[rr]^-{g} \ar[d]_-{n} && B \ar[d]^-{m} \\
A \ar[rr]_-{f} && C}
\end{equation}
in $\CA$ with $m\in \CM$, the square
\begin{equation}
\xymatrix{
P \ar[rr]^-{g_*}  && B  \\
A \ar[rr]_-{f_*} \ar[u]^-{n^*} && C \ar[u]_-{m^*}}
\end{equation}
commutes in $\CP$. 

There is analogue, for partial maps in place of general spans, of the result in
\cite{Lindner1976} about Mackey functors.
There is an equivalence between functors $T \dd \CP \lra \CX$
and pairs $(T^*,T_*)$ of functors 
$T^* \dd \CM^{\mathrm{op}} \lra \CX$
and $T_* \dd \CA \lra \CX$
which agree on objects and satisfy the B\'enabou-Roubaud-Chevalley-Beck condition on pullbacks of the form \eqref{pbm}. 
Indeed $T^*=T\circ (-)^*$ and $T_*=T\circ (-)_*$. 
  
Let $J\dd \CE \ra \CA$ be the inclusion functor and let $I = (-)_* \dd \CA \ra \CP$.
Both functors are the identity on objects.

Let $\CX$ be any category admitting coproducts and products indexed by
sets of $\CM$-subobjects of any given object of $\CA$.

\begin{lemma}\label{leftKan}
Each functor $F\dd \CE \ra \CX$ has a pointwise left Kan extension
\begin{equation}
\begin{aligned}
\xymatrix{
\CE \ar[rd]_{F}^(0.5){\phantom{a}}="1" \ar[rr]^{J}  && \CA \ar[ld]^{\mathrm{Lan}_JF}_(0.5){\phantom{a}}="2" \ar@{=>}"1";"2"^-{\kappa}
\\
& \CX 
}
\end{aligned}
\end{equation} 
along $J\dd \CE \ra \CA$ defined on objects by:
$$(\mathrm{Lan}_JF)X = \sum_{U\le X}FU \ .$$
For morphisms $f \dd X\ra Y$ in $\CA$, the following square commutes.
\begin{equation}
\begin{aligned}
\xymatrix{
FU \ar[rr]^-{\mathrm{in}_U} \ar[d]_-{Fe} && (\mathrm{Lan}_JF)X \ar[d]^-{(\mathrm{Lan}_JF)f} \\
FfU \ar[rr]_-{\mathrm{in}_{fU}} && (\mathrm{Lan}_JF)Y}
\end{aligned}
\end{equation}
\end{lemma}
\begin{proof}
The inclusion $U\le X \mapsto (U, i_U \dd JU\ra X)$ 
of the discrete category on the set $\{U \dd U\le X\}$ into the comma category
$J/X$ has a left adjoint, taking $(V,f\dd JV\ra X)$ to $fV\le X$, and so is final. 
It follows that the colimit of 
$J/X \stackrel{\mathrm{dom}}\lra \CE \stackrel{F}\lra \CX$
can be calculated by restricting along the inclusion and so is the coproduct displayed.
\end{proof}

\begin{lemma}
Each functor $T\dd \CA \ra \CX$ has a pointwise right Kan extension
\begin{equation}
\begin{aligned}
\xymatrix{
\CA \ar[rd]_{T}^(0.5){\phantom{a}}="1" \ar[rr]^{I}  && \CP \ar[ld]^{\mathrm{Ran}_IT}_(0.5){\phantom{a}}="2" \ar@{<=}"1";"2"^-{\rho}
\\
& \CX 
}
\end{aligned}
\end{equation}  
along $I\dd \CA \ra \CP$ defined on objects by:
$$(\mathrm{Ran}_IT)X = \prod_{U\le X}TU \ .$$
For morphisms $(W,h) \dd X\ra Y$ in $\CP$, the square
\begin{equation}
\begin{aligned}
\xymatrix{
(\mathrm{Ran}_IT)X \ar[rr]^-{\mathrm{pr}_{h^{-1}V}} \ar[d]_-{(\mathrm{Ran}_IT)(W,h)} && Th^{-1}V \ar[d]^-{T\bar{h}} \\
(\mathrm{Ran}_IT)Y \ar[rr]_-{\mathrm{pr}_V} && TV 
}
\end{aligned}
\end{equation}
commutes, where $\bar{h}$ is the pullback of $h$ along $i_V$.
\end{lemma}
\begin{proof}
The functor $U\le X \mapsto (i_U^* \dd X\ra IU, U)$ from the discrete category 
on the set $\{U \dd U\le X\}$ into the comma category
$X/I$ has a right adjoint, taking $((W,h)\dd X\ra IA,A)$ to $W\le X$, and so is initial. 
It follows that the limit of 
$X/I \stackrel{\mathrm{cod}}\lra \CA \stackrel{T}\lra \CX$
can be calculated as the product displayed.
\end{proof}

This gives the two adjunctions
\begin{equation}\label{twoadj} 
\xymatrix @R-3mm {
[\CE,\CX] \ar@<1.5ex>[rr]^{\mathrm{Lan}_J} \ar@{}[rr]|-{\bot} && [\CA,\CX]
\ar@<1.5ex>[rr]^{\mathrm{Ran}_I}  \ar@{}[rr]|-{\top} \ar@<1.5ex>[ll]^{[J,1]} && [\CP,\CX]
\ar@<1.5ex>[ll]^{[I,1]}  \ .
}
\end{equation}
Each adjunction generates a comonad on $[\CA,\CX]$:
\begin{equation}\label{twocmnd}
G = \mathrm{Lan}_J \circ [J,1] \ \text{ and } \  H = [I,1] \circ \mathrm{Ran}_I \ .
\end{equation}
To be more explicit, the functors $GF$ and $HF$ are defined on objects by 
$$(GF)A=\sum_{U\le A}{FU} \ \text{ and } \ (HF)A=\prod_{U\le A}{FU} \ .$$ 
On morphisms $f\dd A\ra B$ in $\CA$, they are defined by commutativity in the squares
\begin{equation}
\begin{aligned}
\xymatrix{
FU \ar[rr]^-{Fe} \ar[d]_-{\mathrm{in}_U} && FfU \ar[d]^-{\mathrm{in}_{fU}} \\
(GF)A \ar[rr]_-{(GF)f} && (GF)B}
\qquad
\xymatrix{
(HF)A \ar[rr]^-{(HF)f} \ar[d]_-{\mathrm{pr}_{f^{-1}V}} && (HF)B \ar[d]^-{\mathrm{pr}_V} \\
Ff^{-1}V \ar[rr]_-{F\bar{f}} && FV \ \ .}
\end{aligned}
\end{equation}
The counits $\varepsilon_F \dd GF \ra F$ and $\varepsilon_F \dd HF \ra F$ have
components respectively defined by $\varepsilon_FA \circ \mathrm{in}_U = F\mathrm{in}_U$
and $\varepsilon_FA = \mathrm{pr}_A$. 
The comultiplications $\delta_F \dd GF\ra G^2F$ and $\delta_F \dd HF\ra H^2F$ 
have components respectively defined as follows.
\begin{equation*}
\xymatrix{
FU \ar[rd]_{\mathrm{in}_{U\le U}}\ar[rr]^{\mathrm{in}_U}   && \sum_{U\le A}{FU} \ar[ld]^{\delta_FA} \\
& \sum_{V\le U\le A}{FV}  &
}
\qquad
\xymatrix{
\prod_{U\le A}{FU} \ar[rd]_{\delta_FA}\ar[rr]^{\mathrm{pr}_{V}}   && FV  \\
& \prod_{V\le U\le A}{FV}\ar[ru]_{\mathrm{pr}_{V\le U}}  &
}
\end{equation*}


We shall construct a comonad morphism $\Theta \dd G \Rightarrow H$ when $\CX$ is a pointed
category.

Take $F : \CA \lra \CX$ and $A\in \CA$. We define 
\begin{equation}\label{Thetacmpnt}
\Theta_FA : \sum_{U\le A}FU \lra \prod_{V\le A}FV
\end{equation}
by taking the composite
$$ FU\stackrel{\mathrm{in}_U}\lra \sum_{U\le A}{FU} \stackrel{\Theta_FA}\lra \prod_{V\le A}{FV} \stackrel{\mathrm{pr}_V}\lra FV$$
to be zero unless $U\le V$, in which case the composite is $F(U\le V) \dd FU \ra FV$.

\begin{proposition}
For the two comonads $G$ and $H$ of \eqref{twocmnd} on $[\CA,\CX]$, a comonad morphism $\Theta : G \Rightarrow H$ is defined by the components \eqref{Thetacmpnt}.
\end{proposition}
\begin{proof}
Naturality in $A$ is proved by contemplating the following diagram.
\begin{equation*}
\xymatrix{
FU \ar[d]_{Fe}\ar[r]^{\mathrm{in}_U \ }  & \sum_{U\le A}{FU}\ar[d]^{(GF)f} \ar[r]^{\Theta_FA} & \prod_{U_1\le A}{FU_1} \ar[d]^{(HF)f} \ar[r]^{ \ \mathrm{pr}_{f^{-1}V}} 
& Ff^{-1}V \ar[d]^{F\bar{f}} \\
FfU\ar[r]_{\mathrm{in}_{fU} \ }  & \sum_{V_1\le B}{FV_1}\ar[r]_{\Theta_FB} & \prod_{V\le B}{FV}\ar[r]_{ \ \mathrm{pr}_V} 
& FV }
\end{equation*}
The top composite is zero unless $U\le f^{-1}V$ (say with inclusion $i$). 
The bottom composite is zero unless $fU\le V$ (say with inclusion $j$). 
The conditions are the same.
When this condition holds, the diagram commutes since $\bar{f} \circ i = j \circ e$.  

Naturality in $F$ is obvious. 

Preservation of the counits is proved by the calculation
$$\varepsilon_FA\circ \Theta_FA \circ \mathrm{in}_U = \mathrm{pr}_A\circ \Theta_FA \circ \mathrm{in}_U
= F\mathrm{in}_U = \varepsilon_FA\circ \mathrm{in}_U \ .$$ 

Preservation of the comultiplications is proved by observing that 
$$\mathrm{pr}_{V\le U} \circ \delta_FA \circ \Theta_FA \circ \mathrm{in}_W
= \mathrm{pr}_{V} \circ \Theta_FA \circ \mathrm{in}_W$$
while
$$\mathrm{pr}_{V\le U} \circ \Theta^2_FA \delta_FA  \circ \mathrm{in}_W
= \mathrm{pr}_{V\le U} \circ \Theta^2_FA \circ \mathrm{in}_{W\le W} \ .$$
Both of these right-hand sides are zero unless $W\le V$, in which case they
are both equal to $F(W\le V)$. 
\end{proof}

\begin{proposition}
The functor $[I,1] : [\CP,\CX]\lra [\CA,\CX]$ is comonadic.
\end{proposition}
\begin{proof}
The functor $[I,1]$ is conservative since $I$ is bijective on objects.
It also preserves all limits, including all equalizers.
The result follows (for example) by the Beck comonadicity theorem \cite{ML1963}.
\end{proof}

It follows that $\Theta \dd G\Rightarrow H$
induces a functor $\bar{\Theta} \dd [\CA,\CX]^G \ra [\CP,\CX]$ over $[\CA,\CX]$, 
where $[\CA,\CX]^G$ is the category of Eilenberg-Moore $G$-coalgebras.  
The composite of $\bar{\Theta}$ with the comparison functor 
$[\CE,\CX] \lra [\CA,\CX]^G$ is isomorphic to  
\begin{eqnarray}\label{hatX}
\widehat{M}(F) \dd [\CE, \CX] \lra [\CP, \CX]
\end{eqnarray}
as obtained in Theorem~\ref{combadj}.

Using the finite well-poweredness assumption on the factorization system of $\CA$, 
we shall show that the functor \eqref{hatX} is an equivalence.
In other words, we shall show that the free additive categories on
$\CE$ and $\CP$ are Morita equivalent (= additively Cauchy equivalent).

\begin{theorem}\label{parmapthm}
In the situation described in Example~\ref{Ex1+} and for any additive category $\CX$
with finite direct sums and splitting idempotents, 
the functor \eqref{hatX} is an equivalence of categories
$$[\CE, \CX] \simeq [\CP, \CX] \ . $$
\end{theorem}

We first prove that the comonads $G$ and $H$ of \eqref{twocmnd} are isomorphic.

\begin{proposition}\label{invTheta}
In the situation of Example~\ref{Ex1+} and $\CX$ as in Theorem~\ref{parmapthm}, 
the comonad morphism
$\Theta \dd G\Rightarrow H$ defined at \eqref{Thetacmpnt} is invertible.
\end{proposition}
\begin{proof}
The ordered set of $\CM$-subobjects of $A \in \CA$ is assumed finite
and so has a linear refinement. 
So we can list all the subobjects as $0=U_0, U_1, \dots , U_n = A$ 
such that $U_i \le U_j$ implies $i \le j$.
Then the morphisms \eqref{Thetacmpnt} are represented by an
upper-triangular matrix with identity morphisms in the main diagonal.
Since $\CX$ is additive (including existence of additive inverses), such
matrices are invertible.
\end{proof}

\begin{proof}[Proof of Theorem~\ref{parmapthm}]
By Proposition~\ref{invTheta}, the functor \eqref{hatX} becomes the
comparison from $[\CE,\CX]$ into the category of $G$-coalgebras. 
So the Theorem is now equivalent to comonadicity of the functor 
$\mathrm{Lan}_J \dd [\CE,\CX]\ra [\CA,\CX]$.
Since the unit of the adjunction $\mathrm{Lan}_J \dashv [J,1]$ is a
pointwise coretraction $FA \lra FA\oplus \bigoplus_{U< A}{FU}$, 
and hence a strong monomorphism, 
the functor $\mathrm{Lan}_J$ is conservative. 
So it remains to prove that $\mathrm{Lan}_J$ preserves certain equalizers.
In fact, it preserves all finite limits. Since limits in $[\CA,\CX]$ are formed
pointwise, it suffices to see that each 
$$\mathrm{ev}_A\circ \mathrm{Lan}_J \dd [\CE,\CX]\lra \CX$$
preserves finite limits where $\mathrm{ev}_A \dd [\CA,\CX]\ra \CX$ is evaluation
at $A\in \CA$. By Lemma~\ref{leftKan}, 
$$\mathrm{ev}_A\circ \mathrm{Lan}_J \cong \bigoplus_{U\le A}\mathrm{ev}_U \ .$$
Each evaluation $\mathrm{ev}_U$ preserves limits so their direct sum does.    
\end{proof}

\begin{theorem}\label{settingmainthm}
For the setting of Section~\ref{setting}, the adjunction \eqref{settingkeradj} is an equivalence.
Indeed, the category $\mathrm{Agp}$ of abelian groups can be replaced by any
additive category $\CX$ in which idempotents split and finite direct sums exist,
yielding an equivalence as announced at \eqref{basicequ}. 
\end{theorem}

\section{Examples of Theorem~\ref{settingmainthm}}\label{exx}

\begin{example}\label{babyDPK} 
We begin with a baby version of the Dold-Puppe-Kan Theorem.
Let $\mathrm{Pt}\CX$ denote the category whose objects are split epimorphisms in $\CX$, the morphisms are morphisms of the epimorphisms
which commute with the splittings; this is what Bourn \cite{Bourn91} calls the category of points in $\CX$.  
Take $\mathbb{P}$ to be the free-living adjunction $\mu^*\dashv \mu$ with identity counit $\mu^*\circ \mu = 1$.
So $[\CP,\CX] \cong \mathrm{Pt}\CX$.  
Let $\CM$ consist of all the monomorphisms. 
Then $\CR$ contains only the identities. 
Theorem~\ref{mainthm} yields $$\mathrm{Pt}\CX \simeq \CX \times \CX \ .$$

This example also shows the necessity of $\CX$ having homs enriched in abelian groups (not merely commutative monoids). 
We need $\CX$ to have kernels of split epimorphisms and coproducts already. 
If we also ask that it have finite products then considering the split epimorphism $X\times Y\ra Y$ given by the projection, the counit 
is the canonical map $X+Y \ra X\times Y$, 
so if this is invertible we have hom enrichment in commutative monoids. 
Now considering the codiagonal $X+ X\ra X$, split by one of the 
injections, it is not hard to show that $1_X$ has an additive inverse. 
\end{example}

\begin{example}\label{DPK} (\cite{Dold, DoldPuppe, Kan}) Applying Theorem~\ref{settingmainthm} to 
Example~\ref{Ex2+} yields that $[\Delta_{\bot,\top},\CX]$ is equivalent to 
the category of chain complexes in $\CX$.
\end{example}

\begin{example} 
Applying Theorem~\ref{settingmainthm} to Example~\ref{Ex3+} yields 
$[\mathbb{I},\CX] \simeq [\Delta_{\mathrm{inj}},\CX]$, the category of semi-simplicial objects in $\CX$.
\end{example}
\bigskip

Here are some examples of the general type described in Example~\ref{Ex1+}.
\bigskip

\begin{example} A {\em (set) species} in the sense of Joyal \cite{Species} is a functor 
$F : \mathfrak{S} \lra \mathrm{Set}$
where $\mathfrak{S}$ is the category of sets and bijective functions. 
A {\em pointed-set species} is a functor 
$F : \mathfrak{S} \lra 1/\mathrm{Set}$.
An {\em $R$-module species} is a functor 
$F : \mathfrak{S} \lra \mathrm{Mod}^R$; the case where $R$ is
a field is the basic situation of \cite{AnalFunct}.

Following \cite{CEF2012}, we write $\mathrm{FI}$ for the category of finite sets
and injective functions. We then see that 
$\CM = \mathrm{FI}$ and $\CE = \mathfrak{S}$,
while $\CP = \mathrm{FI}\sharp$, the category of injective partial functions.  

\begin{corollary} \cite{CEF2012} The functor
$$\widehat{M} : [\mathfrak{S}, \mathrm{Mod}^R] \lra [\mathrm{FI}\sharp, \mathrm{Mod}^R]$$
is an equivalence of categories.
\end{corollary}
\end{example}

\begin{example}
Another example relevant to \cite{repfgl} is the category $\CA = \mathrm{FIVect}_{\mathbb{F}_q}$
of finite vector spaces over the field $\mathbb{F}_q$ of cardinality 
$q$ (a prime power) and injective linear functions.
Let $\CE = \mathfrak{Gl}_q$ be the category of finite
$\mathbb{F}_q$-vector spaces and linear isomorphisms. 
Then $\CP = \mathrm{FI}\sharp_q$ is the
category of finite $\mathbb{F}_q$-vector spaces and injective 
partial linear functions.

\begin{corollary} The functor
$$\widehat{M} : [\mathfrak{Gl}_q, \mathrm{Mod}^R] \lra [\mathrm{FI}\sharp_q, \mathrm{Mod}^R]$$
is an equivalence of categories.
\end{corollary}
\end{example}   

\begin{example}
Here are a few examples of categories $\CA$ as in Example~\ref{Ex1+} to which 
Theorem~\ref{settingmainthm} applies with $\CE$ 
the epimorphisms and $\CM$ the monomorphisms:
\begin{enumerate}[(a)]
\item the category of finite abelian groups and group morphisms;
\item the category of finite abelian $p$-groups and group morphisms;
\item the category of finite sets and all functions.
\end{enumerate}
Theorem~\ref{settingmainthm} also applies to $\CM$ in place of $\CA$
in these examples. Then $\CE$ is replaced by the groupoid of
invertible morphisms in $\CA$. In case of example (a),
the paper \cite{HillarDarren2007} describes the groupoid being represented in $\CX$. 
\end{example}

\begin{example}
Consider the ``algebraic'' simplicial category $\Delta_+$ 
whose objects are all the natural numbers and 
whose morphisms $\xi : m\lra n$ are order-preserving functions 
$$\xi : \{0, 1, \dots , m-1\}\lra \{0, 1, \dots , n-1\} \ .$$
Put $\CA = \Delta_+^{\mathrm{op}}$.
Take $\CM$ in $\CA$ to consist of the surjections in $\Delta_+$.
Pushouts of surjections along arbitrary morphisms exist in 
$\Delta_+$. 
Then $\CE = \Delta_{+\mathrm{inj}}^{\mathrm{op}}$
and $\CP$ is the opposite of the category whose morphisms
$m\lra n$ are cospans $$m\stackrel{\xi}\lra r \stackrel{\sigma}\lla n$$
in $\Delta_+$ with $\sigma$ surjective. 
We could also take the ``topological'' simplicial category 
$\Delta$ (omit the object $0$) to obtain
a reinterpretation of the preoperads in $\CX$ in the sense of 
\cite{Berger1996}.
\end{example}

\begin{example}
Here is a rather trivial example involving $\Delta$.
Take $\CA$ to be the category of non-empty ordinals and morphisms
the order-preserving functions which preserve first element.
Let $\CM$ be the class of morphisms which are inclusions of
initial segments. This is part of a factorization system where
$\CE = \Delta_{\bot \neq \top}$ is the category of ordinals with 
distinct first and last element 
and morphisms the order-preserving functions which preserve first 
and last element. Sometimes $\CE$ is called the category of {\em intervals}; there is a duality isomorphism
$$\CE \cong \Delta^{\mathrm{op}} \ .$$ 
In this case, not only do we have the equivalence
$$[\CE, \CX] \simeq [\CP, \CX] $$
of Theorem~\ref{settingmainthm}, we actually also have an isomorphism
$$\CP \cong \CE \ .$$
\end{example}

\begin{example}
Take $\CA$ to be a (partially) ordered set with finite infima and the descending chain condition. 
Then every morphism is a monomorphism
and the strong epimorphisms are equalities. So $\CE$ is
the discrete category $\mathrm{ob}\CA$ on the set of
elements of the ordered set. The reader may like to contemplate
the case where $\CA$ is the set of strictly positive integers
ordered by division. 
\end{example}

\section{Remarks on idempotents}\label{idemp}

Define a relation on any monoid $M$ by $a\sqsubseteq b$ when $ba = a$.
This relation is transitive; indeed, we have a stronger property.

\begin{proposition}\label{gentrans}
If $ub\sqsubseteq a$ and $vc \sqsubseteq b$ then $uvc\sqsubseteq a$.

If $u_ia_i\sqsubseteq a_{i-1}$ for $1\sqsubseteq i\sqsubseteq n$ then $u_1\dots u_na_n\sqsubseteq a_0$.
\end{proposition}
\begin{proof}
For the first sentence, the assumptions are $aub = ub$ and $bvc=vc$.
So $auvc = aubvc=ubvc=uvc$; that is, $uvc \sqsubseteq a$. 
The second sentence follows by induction.
\end{proof}

The relation is only reflexive for idempotents: clearly $a\sqsubseteq a$ is equivalent to $aa=a$.

The unit $1$ of the monoid is a largest element in the sense that $a\sqsubseteq 1$ for all $a\in M$.
That is, $1$ is the empty meet. 
However, not all meets need exist. A {\em meet} for $a,b\in M$ is an element $a\wedge b$
with $a\wedge b \sqsubseteq a$ and $a\wedge b \sqsubseteq b$, and, if $x\sqsubseteq a$ and $x\sqsubseteq b$ then
$x\sqsubseteq a\wedge b$. In particular, $a\wedge b$ must be an idempotent. 
Meets of lists of $n$ elements are defined in the obvious way and, for $n\ge 2$,
can be constructed from iterated binary meets when they exist. 

\begin{proposition}\label{someets}
If $ab\sqsubseteq b$ and $a\sqsubseteq a$ then $a\wedge b = ab$.

If $a_1, \dots , a_n$ are idempotents such that $a_ia_j\sqsubseteq a_j$ for $i\sqsubseteq j$
then $a_1\wedge \dots \wedge a_n = a_1 \dots a_n$. 
\end{proposition}
\begin{proof}
For the first sentence, we are told that $ab\sqsubseteq b$, while $a\sqsubseteq a$ implies $aa=a$,
and so $aaab=ab$, yielding $ab\sqsubseteq a$.
For the second sentence the result is clear for $n=1$ since we suppose $a_1$
idempotent. Assume the result for $n-1$; so $a_1\wedge \dots \wedge a_{n-1} = a_1 \dots a_{n-1}$. Apply the second sentence of Proposition~\ref{gentrans} to the inequalities $a_ia_n\sqsubseteq a_n$ to deduce $a_1\dots a_{n-1}a_n\sqsubseteq a_n$.
So, by the first sentence, $a_1\dots a_n = a_1\dots a_{n-1}\wedge a_n = a_1\wedge \dots \wedge a_{n-1}\wedge a_n$, as required.
\end{proof}

Notice that, if $ab=ba$ and $b$ is idempotent, then $bab=abb=ab$, so $ab\sqsubseteq b$.
So the proposition applies to commuting idempotents.

Now suppose we have a ring $R$. We can apply our results to the multiplicative monoid of $R$. 
We say idempotents $e$ and $f$ in $R$ are {\em orthogonal}
when $ef=fe=0$. A list $e_0, e_1, \dots e_n$ of idempotents is {\em orthogonal}
when each pair in the list is orthogonal. The list is {\em complete} when
$e_0+ e_1+ \dots + e_n = 1$. An easy induction shows that a complete
list of idempotents is orthogonal if and only if $e_ie_j = 0$ for $i\le j$.
 
For each $a\in R$, put $\bar{a}=1-a$. Clearly if $a$ is idempotent, so is $\bar{a}$.

Let $R^{\circ}$ denote the ring obtained from $R$ by reversing multiplication.

\begin{proposition}
\begin{enumerate}
\item[(a)] $a\sqsubseteq b$ in $R$ if and only if $\bar{b}\sqsubseteq \bar{a}$ in $R^{\circ}$. 
\item[(b)] For $b$ an idempotent, 
$ab\sqsubseteq b$ in $R$ if and only if $\bar{a}\bar{b}\sqsubseteq \bar{b}$ in $R^{\circ}$.
\item[(c)] If $a$ and $b$ are idempotents and $ab\sqsubseteq b$ in $R$ then 
$e_0=ab, e_1=\bar{a}b, e_2=\bar{b}$
is a complete list of orthogonal idempotents. 
\end{enumerate}
\end{proposition}
\begin{proof}
\begin{enumerate}
\item[(a)] $\bar{b}\sqsubseteq \bar{a}$ in $R^{\circ}$ means $(1-b)(1-a) = 1-b$ in $R$; 
that is, $ba=a$ which means $a\sqsubseteq b$ in $R$. 
\item[(b)] $\bar{a}\bar{b}\sqsubseteq \bar{b}$ in $R^{\circ}$ means 
$\bar{b}\bar{a}\bar{b}= \bar{b}\bar{a}$ in $R$. That is, 
$(1-b)(1-a)(1-b) = (1-b)(1-a)$.
That is, $1-a-b+ab-b+ba+bb-bab = 1-b-a+ba$.
That is, $bab = ab$, which is $ab\sqsubseteq b$ in $R$. 
\item[(c)] We already know $e_0$ and $e_2$ are idempotent.
They are also orthogonal: $e_0e_2 = ab(1-b)=ab-ab=0$
and $e_2e_0=(1-b)ab=ab-bab=0$. 
Therefore $e_0+e_2 = ab + 1 - b = 1-(1-a)b = \overline{e_1}$
is idempotent. So $e_1$ is idempotent and $e_0+e_1+e_2=1$.
The calculations $e_0e_1=ab(1-a)b=ab-abab=0$ and 
$e_1e_2=(1-a)b(1-b)=b-ab-b+abb=0$ complete the proof.
\end{enumerate}
\end{proof}

We can extend part (c) inductively to obtain: 

\begin{proposition}\label{weakcommute}
Suppose $a_1, \dots , a_n$ are idempotents such that $a_ia_j\sqsubseteq a_j$ for $i\le j$.
Then $e_i = \overline{a_i}a_{i+1}a_{i+2} \dots a_{n}$ for $0\le i \le n$
(in particular, $e_0 = a_1a_2\dots a_n$ and $e_n = \overline{a_n}$)
defines a complete list of orthogonal idempotents. 
\end{proposition}

Suppose $\CX$ is an additive category in which idempotents split.
Our results apply to the endomorphism monoid $\CX (A,A)$ of
each object $A\in \CX$. If $a$ is an idempotent on $A$, we have
a splitting:
\begin{equation*}
\xymatrix{
A \ar[rr]^-{a} \ar[d]_-{r_a} && A \ar[d]^-{r_a} \\
aA \ar[rr]_-{1_{aA}} \ar[rru]_-{i_{a}} && aA \ .}
\end{equation*} 
Yet, we also have a splitting for $\bar{a} = 1-a$ which incidentally provides a kernel
$\bar{a}A$ for $a$ and so a direct sum decomposition of $A$:
$$A \cong \bar{a}A\oplus aA \ .$$

More generally, for any complete list $e_0, e_1, \dots e_n$ of orthogonal idempotents
in $\CX (A,A)$, we obtain a direct sum decomposition
$$A\cong e_0A \oplus e_1A \oplus \dots \oplus e_nA \ .$$ 

\section{When $\CX$ is semiabelian}\label{semiabelian}

Semiabelian categories include the category $\mathrm{Grp}$ of 
(not necessarily abelian) groups and group morphisms.
In \cite{Bourn07} Dominique Bourn gave a version of the Dold-Puppe-Kan Theorem (Example~\ref{DPK})
for the case where the codomain category $\CX$ was semiabelian.
In that case it asserted monadicity of the right adjoint in 
\begin{equation}\label{pttedkeradj} 
\xymatrix @R-3mm {
[\CD,\CX]_{\mathrm{pt}} \ar@<1.5ex>[rr]^{\widehat{M}} \ar@{}[rr]|-{\bot} && [\CP,\CX]
 \ar@<1.5ex>[ll]^{\widetilde{M}} \ .
}
\end{equation} 
In this section, we provide a version of this for $\CP$ as in Section~\ref{setting}
in place of $\Delta^{\mathrm{op}}$. However, we require an extra assumption related to
idempotents. 

\begin{assumption}\label{ass6}
The maximal proper elements of $\mathrm{Sub}A$ can be listed $m_1, \dots, m_n$ 
such that the idempotents $c_i = m_i \circ m_i^*$ on $A$ satisfy 
$c_j\circ c_i\circ c_j=c_j\circ c_i$ for all $i< j$. 
\end{assumption}

Throughout we assume our category $\CX$ has zero morphisms (that is, has homs enriched in pointed sets).

We begin by providing a non-additive version of the material at the end of Section~\ref{idemp} on idempotents. 

\begin{proposition}
Suppose the category $\CX$ has kernels of idempotents.
Let $e,f$ be idempotents on an object $A$ of $\CX$.
If $e\circ f\circ e = e\circ f$ then the intersection of the kernels of $e$ and $f$ exists.
\end{proposition}
\begin{proof}
Let $k:K\ra A$ be the kernel of $e$.
Since $e\circ f \circ k=e\circ f\circ e\circ f =0$, there exists a unique $g$ with 
$f\circ k=k\circ g$. Then $k\circ g\circ g = f\circ k\circ g=f\circ f \circ k=f\circ k=k\circ g$ and $k$ is a monomorphism. So $g$ is idempotent. Then the kernel $\ell :L\ra K$ of $g$ is easily verified to be the intersection of the kernels of $e$ and $f$.
\end{proof}

Protomodular categories were defined by Bourn \cite{Bourn91}:
a category $\CX$ (with zero morphisms) is {\em protomodular} when 
it is finitely complete and, for each object $A$, 
the functor $\mathrm{ker}: \mathrm{Pt}A\ra \CX$ is conservative. 
Here $\mathrm{Pt}A$ is the category
whose objects $(p,X,s)$ consist of morphisms $p:X\ra A, s:A\ra X$ with
$p\circ s = 1_A$, and whose morphisms $f:(p,X,s)\ra (q,Y,t)$ are morphisms
$f:X\ra Y$ such that $q\circ f = p$ and $f\circ s = t$.
Also the functor $\mathrm{ker}$ takes $(p,X,s)$ to the kernel of $p$.

The following property is sometimes \cite{BorBou07} taken as the definition of protomodular.  

\begin{lemma}\label{protolemma} 
In a protomodular category, if $(p,X,s)$ is an object of 
$\mathrm{Pt}A$ and $k:K\ra X$
is the kernel of $p$ then $s:A\ra X, k:K\ra X$ are jointly strongly epimorphic.
\end{lemma}
\begin{proof} Suppose $m:Y\ra X$ is a monomorphism and $m\circ u = s, m\circ v = k$ for some $u,v$. Then $m:(p\circ m,Y,u)\ra (p,X,s)$ is a morphism of
$\mathrm{Pt}A$. Using $v$, we see that $m$ induces an isomorphism between
the kernel of $p\circ m$ and $K$. Since $\mathrm{ker}: \mathrm{Pt}A\ra \CX$
is conservative, $m:(p\circ m,Y,u)\ra (p,X,s)$ is invertible. So $m$ is invertible.   
\end{proof}      

\begin{proposition}\label{protoidemp}
Let $a_1,\dots,a_n$ be a list of idempotents on an object $A$ of a protomodular category $\CX$. 
Suppose $a_i\circ a_j\circ a_i=a_i\circ a_j$ for $i<j$.
Suppose $a_i=m_i\circ m_i^*$ is a splitting of $a_i$ via a subobject $m_i:A_i\ra A$
and retraction $m_i^*$.
Let $k_i:K_i\ra A$ be the kernel of $m_i^*$ (or equally of $a_i$).
Then the morphisms $m_1,\dots , m_n$ along with the inclusion $\cap_i{\mathrm{ ker} \ m_i^*} \ra A$ are jointly strongly epimorphic.
\end{proposition}
\begin{proof}
By Lemma~\ref{protolemma}, for each $i$, the morphisms $m_i:A_i\ra A$
and $k_i:K_i\ra A$ are jointly strongly epimorphic; we will loosely say ``$A_i$ and $K_i$
cover $A$''. 

If $i>1$ then $a_1\circ a_i\circ k_1=a_1\circ a_i\circ a_1\circ k_1=0$, and so $a_i\circ k_1$
lands in $K_1$, providing a factorization $a_i\circ k_1=k_1\circ a_i^1$.
Now $a_i^1$ is also an idempotent, and, for $1<i<j$, $k_1\circ a_i^1\circ a_j^1\circ a_i^1= a_i\circ a_j\circ a_i\circ k_1 = a_i\circ a_j\circ k_1= k_1\circ a_i^1\circ a_j^1$, and so
$a^1_i\circ a^1_j\circ a^1_i= a^1_i\circ a^1_j$. Clearly the splitting $A^1_i$ of $a^1_i$
is contained in $A_i$.

The kernel $K^1_i$ of $a^1_i$ is $K^1_i=K_1\cap K_i$ since $a_1\circ x=a_i\circ x = 0$ is equivalent to
$x=k_1\circ y$ and $k_1\circ a^1_i\circ y = a_i\circ k_1\circ y = s_i\circ x = 0$; 
so in fact $a^1_i\circ y=0$. 

We know that $A$ may be covered by $A_1$ and $K_1$. By Lemma~\ref{protolemma} again, we know, for each $i>1$, that $K_1$ may be covered by the splitting of $A^1_i$
and the kernel $K^1_i=K_1\cap K_i$ of $a^1_i$. Since $A^1_i\le A_i$, we see that
$A$ may be covered by $A_1$, $A_i$, and $K_1\cap K_i$.

Now continue inductively.  
\end{proof}

\begin{theorem}\label{protocounit}
Suppose $\CP$ is as in Section~\ref{setting} and $\CX$ is protomodular (with zero morphisms) with finite coproducts.
Then the components of the adjunction \eqref{pttedkeradj} are strong epimorphisms.
\end{theorem}
\begin{proof}
The counit has components
$$\sum_{B\preceq_mA}{\bigcap_{C\prec_n B}{\mathrm{ker} \ Tn^*}} \lra TA \ .$$
We prove this is a strong epimorphism by induction on the number $k$ of maximal
proper subobjects $A_1, \dots, A_k$ of $A$, with $m_i:A_i\ra A$. 
The result is clear for $k=1$. For $k>1$, consider the following diagram.
\begin{eqnarray*}
\xymatrix{
\sum_i{\sum_{B\preceq_mA}{\bigcap_{C\prec_n B}{\mathrm{ker} \ Tn^*}}} \ar[rr]^-{ } \ar[d]_-{\delta} && \sum_i{TA_i} \ar[d]^-{\gamma} && \\
\sum_{B\prec_mA}{\bigcap_{C\prec_n B}{\mathrm{ker} \ Tn^*}} \ar[rr]_-{\alpha} && FA && \bigcap_{C\prec_n A}{{\mathrm{ker} \ Tn^*}} \ar[ll]^-{\beta}}
\end{eqnarray*}
The component of the counit is strongly epimorphic if and only if $\alpha$ and
$\beta$ are jointly strongly epimorphic.
The top row is a coproduct of components of the counit already known to be strongly
epimorphic by induction. 
So it suffices to show that $\gamma$ and $\beta$ are jointly strongly epimorphic.
Rewriting the domain of $\beta$ as
$$\bigcap_{C\prec_n A}{{\mathrm{ker} \ Tn^*}} = \bigcap_{i=1}^k{{\mathrm{ker} \ Tm_i^*}} \ ,$$
we see that Proposition~\ref{protoidemp} applies to yield what we want.
\end{proof}

A category $\CX$ is {\em semiabelian} \cite{JMT} when it has zero morphisms, 
is protomodular, is Barr exact, and has finite coproducts.
A category is {\em regular} \cite{Excat} when it is finitely complete, and has the 
(strong epimorphism, monomorphism)-factorization system existing 
and stable under pullbacks. 
It follows that every strong epimorphism is regular (that is, a coequalizer);
see \cite{CarSt86} for a proof. 
A category is {\em Barr exact} when it is regular and every equivalence relation 
is a kernel pair. 

We mentioned Bourn's category $\mathrm{Pt} \CX$ in Example~\ref{babyDPK}.
We will use the following routine fact.

\begin{lemma}\label{ptpresstepi}
If $\CX$ is a semiabelian category then the functor 
$\mathrm{Pt} \CX\ra \CX$, sending each split epimorphism to its kernel, 
preserves strong epimorphisms.
\end{lemma}
\begin{proof}
A strong epimorphism
\begin{eqnarray*}
\xymatrix{
X \ar[rr]^-{g} \ar@<1ex>[d]^{p} && Y \ar@<1ex>[d]^{q} \\
A\ar@<1ex>[u]^{s} \ar[rr]_-{f} && B\ar@<1ex>[u]^{t}}
\end{eqnarray*}
in $\mathrm{Pt} \CX$ has $f$ and $g$ strong epimorphisms in $\CX$.
From this it is easily verified that the square involving the downward-pointing
arrows is a pushout. 
Factor the morphism in $\mathrm{Pt} \CX$ as 
\begin{eqnarray*}
\xymatrix{
X \ar[rr]^-{g_1} \ar@<1ex>[d]^{p} && Z\ar[rr]^-{f_1} \ar@<1ex>[d]^{q_1} && Y \ar@<1ex>[d]^{q} \\
A\ar@<1ex>[u]^{s} \ar[rr]_-{1_A} && A\ar[rr]_-{f} \ar@<1ex>[u]^{t_1} && B\ar@<1ex>[u]^{t}
}
\end{eqnarray*}
in which the right-hand square involving the downward-pointing arrows is a pullback.
The induced morphism $\mathrm{ker} q_1\ra \mathrm{ker} q$ is invertible.
Semiabelian categories are Maltsev \cite{JMT}. 
Therefore, as a comparison morphism  to
the pullback in a pushout square in a Maltsev exact category, $g_1$ is a strong epimorphism (see Theorem 5.7 of \cite{CKP93}). 
The induced morphism $\mathrm{ker} p\ra \mathrm{ker} q_1$ is the pullback
of $g_1$ along $\mathrm{ker} q_1 \ra Z$ and so is a strong epimorphism by regularity.
\end{proof}

\begin{theorem}
If $\CP$ is as in Section~\ref{setting} and $\CX$ is semiabelian then the adjunction
of Theorem~\ref{adjthm} is (crudely) monadic.
\end{theorem}
 \begin{proof}
 By Theorem~\ref{protocounit}, the right adjoint (tilde) is conservative (since this is 
 logically equivalent to the counit being a strong epimorphism).
 Since $\CX$ is semiabelian, it has coequalizers. 
 Therefore $[\CP,\CX]$ has coequalizers, so, for crude monadicity \cite{CWM}, 
 it suffices to show that tilde preserves coequalizers of reflexive pairs.
 
 Both $[\CP,\CX]$ and $[\CD,\CX]_{\mathrm{pt}}$ are semiabelian.
 A limit-preserving functor between semiabelian categories preserves coequalizers of
 reflexive pairs provided it preserves strong (= regular) epimorphisms 
 (see Lemma 5.1.12 of \cite{BorBou04}).
 
 Let $q:S\ra T$ be a strong epimorphism in $[\CP,\CX]$.
 Each $q_A:SA\ra TA$ is a strong epimorphism.
 We must show each $\widetilde{q}_A:\widetilde{S}A\ra \widetilde{T}A$ is a strong epimorphism.
 
 Recall that $\widetilde{T}A$ is calculated by a sequence of kernels of split epimorphisms.
 This sequence depends only on the object $A$ and the category $\CP$, not on the
 particular functor $T$. 
 The desired result follows on repeated application of Lemma~\ref{ptpresstepi}.   
  \end{proof}

\begin{thebibliography}{00}

\bibitem{Excat} Michael Barr, Pierre A. Grillet, Donovan H. van Osdol, \textit{Exact Categories and Categories of Sheaves}, Lecture Notes in Math. \textbf{236} 
(Springer-Verlag, 1971).

\bibitem{UPMT} Michael Batanin and Ross Street, \textit{The universal property of the multitude of trees}, J. Pure Appl. Algebra \textbf{154} (2000) 3--13.\label{UPMT}

\bibitem{Berger1996} Clemens Berger, \textit{Op\'erades cellulaires et espaces de lacets it\'er\'es}, Annales de L'Institut Fourier \textbf{46(4)} (1996) 1125--1157.\label{Berger1996} 

\bibitem{BorBou04} Francis Borceux and Dominique Bourn, \textit{Mal'cev, protomodular, homological and semi-abelian categories}, Mathematics and its Applications \textbf{566} (Kluwer Academic Publishers, Dordrecht, 2004).

\bibitem{BorBou07} Francis Borceux and Dominique Bourn, \textit{Split extension classifier and centrality}, in ``Categories in Algebra, Geometry and Mathematical Physics'', 
Contemporary Mathematics \textbf{431} (2007) 85--104.\label{BorBou07}

\bibitem{Bourn91} Dominique Bourn, \textit{Normalization equivalence, kernel equivalence and affine categories category theory}, Lecture Notes in Math. \textbf{1488} (Springer-Verlag, 1991) 43--62.\label{Bourn91}

\bibitem{Bourn07} Dominique Bourn, \textit{Moore normalization and Dold-Kan theorem for semi-abelian categories}, 
in ``Categories in Algebra, Geometry and Mathematical Physics'', 
Contemporary Mathematics \textbf{431} (2007) 105--124.\label{Bourn07}

\bibitem{CKP93} Aurelio Carboni, G. Max Kelly and M. Cristina Pedicchio, \textit{Some remarks on Maltsev and Goursat categories}, Applied Categorical Structures \textbf{1} (1993) 385--421.\label{CKP93}

\bibitem{CarSt86} Aurelio Carboni and Ross Street, \textit{Order ideals in categories}, 
Pacific J. Math. \textbf{124} (1986) 275--288.\label{CarSt86}

\bibitem{CEF2012} Thomas Church, Jordan S. Ellenberg and Benson Farb, \textit{FI-modules: a new approach to stability for $S_n$-representations} (\url{arXiv:1204.4533v2}, 28 Jun 2012).\label{CEF2012}

\bibitem{CransThesis} Sjoerd Crans, \textit{On combinatorial models for higher dimensional homotopies}, (PhD Thesis, Universiteit Utrecht, 1995).\label{CransThesis}

\bibitem{Dold} Albrecht Dold, \textit{Homology of symmetric products and other functors of complexes}, Annals of Math. \textbf{68} (1958) 54--80.\label{Dold}
 
\bibitem{DoldPuppe} Albrecht Dold and Dieter Puppe, \textit{Homologie nicht-additiver Funktoren. Anwendungen}, Ann. Inst. Fourier Grenoble \textbf{11} (1961) 201--312.\label{DoldPuppe}

\bibitem{FreydKelly} Peter J. Freyd and G. Max Kelly, \textit{Categories of continuous functors I}, J. Pure and Applied Algebra \textbf{2} (1972) 169--191; Erratum Ibid. \textbf{4} (1974) 121.\label{FreydKelly}

\bibitem{GabZis} Peter Gabriel and Michel Zisman, \textit{Calculus of Fractions and Homotopy Theory} (Springer-Verlag, 1967).\label{GabZis} 

\bibitem{Helm} Randall D. Helmstutler, \textit{Model category extensions of the Pirashvili-S?omi?ska theorems} \url{arXiv:0806.1540}.\label{Helm}

\bibitem{HillarDarren2007} Christopher J. Hillar and Darren L. Rhea, \textit{Automorphisms of finite abelian groups}, American Mathematical Monthly \textbf{114(10)} (2007) 917--923; \url{arXiv:math/0605185}.\label{HillarDarren2007}

\bibitem{JMT} George Janelidze, L\'aszl\'o M\'arki, Walter Tholen, \textit{Semi-abelian categories}, J. Pure and Applied Algebra \textbf{168} (2002) 367--386.\label{JMT}

\bibitem{Species} Andr\'e Joyal, \textit{Une theory combinatoire des series formelles}, Advances in Mathematics \textbf{42} (1981) 1--82.\label{Species}

\bibitem{AnalFunct} Andr\'e Joyal, \textit{Foncteurs analytiques et esp\`eces de structures}, Lecture Notes in Mathematics \textbf{1234} (Springer 1986) 126--159.\label{AnalFunct} 

\bibitem{repfgl} Andr\'e Joyal and Ross Street, \textit{The category of representations of the general linear groups over a finite field}, J. Algebra \textbf{176 (3)} (1995) 908--946.\label{repfgl}

\bibitem{Kan} Daniel Kan, \textit{Functors involving c.s.s complexes}, Transactions of the American Mathematical Society \textbf{87} (1958) 330--346.\label{Kan}

\bibitem{KelSt1974} G. Max Kelly and Ross Street, \textit{Review of the elements of 2-categories}, Lecture Notes in Mathematics \textbf{420} (Springer-Verlag, 1974) 75--103.\label{KelSt1974}

\bibitem{KM09} Ganna Kudryavtseva and Volodymyr Mazorchuk, 
\textit{On Kiselman's semigroup}, Yokohama Math. J. \textbf{55} (2009) 21--46.\label{KM09} 

\bibitem{Law91} F. William Lawvere, \textit{More on Graphic Toposes}, Cahiers de topologie et g\'eom\'etrie diff\'erentielle \textbf{32(1)} (1991) 5--10.\label{Law91}

\bibitem{Lindner1976} Harald Lindner, \textit{A remark on Mackey functors}, 
Manuscripta Mathematica \textbf{18} (1976) 273--278.\label{Lindner1976}

\bibitem{ML1963} Saunders Mac Lane, \textit{Natural associativity and commutativity}, 
Rice University Studies \textbf{49} (1963) 28--46.\label{ML1963}

\bibitem{CWM} Saunders Mac Lane, \textit{Categories for the Working Mathematician}, 
Graduate Texts in Mathematics \textbf{5} (Springer-Verlag, 1971).\label{CWM} 

\bibitem{Mur38} Francis D. Murnaghan, \textit{The Analysis of the Kronecker Product of Irreducible Representations of the Symmetric Group}, American J. Math. \textbf{60(3)}(1938) 761--784.\label{Mur38}

\bibitem{Mur55} Francis D. Murnaghan, \textit{On the analysis of the Kronecker product of irreducible representations of $S_n$}, Proc. Nat. Acad. Sci. U.S.A. \textbf{41} (1955) 515--518.\label{Mur55}

\bibitem{PanSt} Elango Panchadcharam and Ross Street, \textit{Mackey functors on compact closed categories}, Journal of Homotopy and Related Structures \textbf{2(2)} (2007) 261--293.\label{PanSt}

\bibitem{Pirash00} Teimuraz Pirashvili, \textit{Dold-Kan type theorem for $\Gamma$-groups}  
Mathematische Annalen \textbf{318} (2000) 277--298.\label{Pirash00}

\bibitem{Segal} Graeme Segal, \textit{Categories and cohomology theories}, 
Topology \textbf{13} (1974) 293--312.\label{Segal}

\bibitem{Slom} Jolanta S\l omi\'nska, \textit{Dold-Kan type theorems and Morita equivalences of functor categories}, Journal of Algebra \textbf{274} (2004)118--137.\label{Slom}

\bibitem{Sta99} Richard P. Stanley, \textit{Enumerative combinatorics, Vol. 2}, Cambridge Studies in Advanced Mathematics \textbf{62} (Cambridge University Press, 1999).\label{Sta99} 

\bibitem{ECC} Ross Street, \textit{Enriched categories and cohomology with author commentary}, Reprints in Theory and Applications of Categories \textbf{14} (2005) 1--18;
originally Quaestiones Math. \textbf{6} (1983) 265--283.\label{ECC}

\bibitem{ACEC} Ross Street, \textit{Absolute colimits in enriched categories}, Cahiers de topologie et g\'eom\'etrie diff\'erentielle \textbf{24} (1983) 377-- 379.\label{ACEC}

\bibitem{LDTHOC} Ross Street, \textit{Low-dimensional topology and higher-order categories}, International Conference on CategoryTheory CT95, Halifax, July 9--15 1995; \url{http://www.mta.ca/~cat-dist/ct95.html} and \url{http://maths.mq.edu.au/~street/LowDTop.pdf}.\label{LDTHOC}

\bibitem{Verity07} Dominic Verity, \textit{Weak complicial sets. II. Nerves of complicial Gray-categories}, 
in ``Categories in Algebra, Geometry and Mathematical Physics'', 
Contemporary Mathematics \textbf{431} (2007) 441--467.\label{Verity07}

\end{thebibliography}

\end{document}